\documentclass[11pt]{article}
\usepackage[T1]{fontenc}
\usepackage{lmodern}
\usepackage[margin=0.92in]{geometry}
\usepackage{amsmath,amssymb,amsthm,mathtools}
\usepackage{booktabs,array,longtable,tabularx}
\usepackage{xcolor,microtype,fancyhdr,enumitem}
\usepackage[hyphens]{url}
\usepackage{xurl}
\usepackage{listings}
\usepackage{graphicx}
\usepackage[hypertexnames=false,colorlinks=true,linkcolor=blue!45!black,
  citecolor=blue!45!black,urlcolor=blue!45!black]{hyperref}
\usepackage[nameinlink,noabbrev]{cleveref}

\setlength{\emergencystretch}{2.5em}
\setlist{itemsep=2pt,topsep=4pt}
\lstset{basicstyle=\ttfamily\footnotesize,breaklines=true,columns=fullflexible,
  keepspaces=true,frame=single,rulecolor=\color{black!20}}
\pagestyle{fancy}\fancyhf{}
\fancyhead[L]{\small Hypergeometric functions in iPRA}
\fancyhead[R]{\small Finite algebra and effective analysis}
\fancyfoot[C]{\thepage}
\setlength{\headheight}{14pt}

\newtheorem{theorem}{Theorem}[section]
\newtheorem{lemma}[theorem]{Lemma}
\newtheorem{proposition}[theorem]{Proposition}
\newtheorem{corollary}[theorem]{Corollary}
\newtheorem{claim}[theorem]{Claim}
\theoremstyle{definition}
\newtheorem{definition}[theorem]{Definition}
\newtheorem{construction}[theorem]{Construction}
\newtheorem{contract}[theorem]{Effective contract}
\newtheorem{example}[theorem]{Example}
\theoremstyle{remark}
\newtheorem{remark}[theorem]{Remark}
\newtheorem{audit}[theorem]{Formalization audit}

\newcommand{\N}{\mathbb N}
\newcommand{\Z}{\mathbb Z}
\newcommand{\Q}{\mathbb Q}
\newcommand{\R}{\mathbb R}
\newcommand{\C}{\mathbb C}
\newcommand{\ipra}{\mathsf{iPRA}}
\newcommand{\PR}{\mathsf{PR}}
\newcommand{\F}[4]{{}_2F_1\!\left(\begin{matrix}#1,#2\\#3\end{matrix};#4\right)}
\newcommand{\PFQ}[4]{{}_{#1}F_{#2}\!\left(\begin{matrix}#3\end{matrix};#4\right)}
\newcommand{\poch}[2]{\left(#1\right)_{#2}}
\newcommand{\abs}[1]{\left|#1\right|}
\newcommand{\norm}[1]{\left\lVert#1\right\rVert_1}
\newcommand{\dd}{\,\mathrm d}
\newcommand{\thetaop}{\vartheta}
\newcommand{\code}[1]{\texttt{#1}}
\newcommand{\StatusA}{\textsf{A}}
\newcommand{\StatusB}{\textsf{B}}
\newcommand{\StatusC}{\textsf{C}}
\newcommand{\StatusD}{\textsf{D}}
\newcommand{\bigO}{\mathcal O}
\DeclareMathOperator{\RePart}{Re}
\DeclareMathOperator{\ImPart}{Im}
\DeclareMathOperator{\Coeff}{Coeff}
\DeclareMathOperator{\Tail}{Tail}
\DeclareMathOperator{\ord}{ord}

\title{Gauss's Hypergeometric Function in Intuitionistic\\
Primitive Recursive Arithmetic}
\author{A constructive proof package prepared for Vladimir Reshetnikov}
\date{16 November 2026}
\hypersetup{pdftitle={Gauss's Hypergeometric Function in Intuitionistic Primitive Recursive Arithmetic},
 pdfauthor={Prepared for Vladimir Reshetnikov},
 pdfsubject={Human proofs, quantitative iPRA reductions, and Lean formalization plan}}

\begin{document}
\maketitle

\begin{abstract}
We reconstruct a large portion of the standard theory of Gauss's
hypergeometric function in a form suitable for intuitionistic primitive
recursive arithmetic.  The central separation is between exact finite
algebra and effective analysis.  The coefficient recurrence, termination,
derivative identities, the hypergeometric differential equation, Frobenius
solutions, Euler--Pfaff transformations, contiguous relations, finite
continued fractions, Chu--Vandermonde, Jacobi and Legendre polynomial
specializations, elementary reductions, and coefficientwise confluent limits
are reduced to primitive recursive finite products, sums, and
quantifier-free induction.  Convergence, differentiation of the infinite
series, Euler's integral, Gauss--Kummer--Bailey summations, connection
formulas, and analytic continuation are stated as first-order theorem schemas
for fixed primitive recursive rational or complex-rational names carrying
explicit bounds, moduli, and apartness margins.

The result is deliberately not a classical assertion quantifying over a
completed field of all real or complex numbers.  Every equality of analytic
objects means a uniform family of rational inequalities with a primitive
recursive precision schedule.  We give complete human arguments for the
stated effective contracts, identify the finite induction invariants, and
explain how to compile them into ProveIt's embedded calculus
\code{Prf PRAAx} and then through its native bi-interpretation.  A companion
Lean file contains executable finite definitions and theorem signatures; it
is an outline, not a claimed kernel-checked analytic formalization.  Exact
rational validation code, a coverage ledger, and source/provenance manifests
are included.
\end{abstract}

\medskip
\noindent\textbf{Status key.}
\StatusA{} = exact finite/formal-power-series theorem in iPRA;
\StatusB{} = effective analytic theorem with explicit primitive recursive
certificates;
\StatusC{} = conditional theorem whose additional contour, branch, or
nonvanishing certificate is displayed;
\StatusD{} = not proved here (usually a global topological or modular
classification claim).

\newpage
\begingroup\small\tableofcontents\endgroup
\newpage

\section{Scope, semantics, and the main result}\label{sec:scope}

\subsection{What is being proved}

The two reference pages mix several logically different kinds of statement:
finite coefficient identities, convergence assertions, differential-equation
claims, integral representations, branch-dependent continuation formulas,
and global assertions about monodromy or modular functions.  Treating all of
these as though they had the same foundational content would conceal the main
issue.  In first-order \(\ipra\), variables range over natural numbers.  There
is no primitive sort of arbitrary real numbers, no quantification over all
complex functions, and no completeness axiom that manufactures limits without
rates.  We therefore use the following four-level interpretation.

\begin{definition}[Four proof levels]\label{def:levels}
A displayed classical statement is assigned one of the following readings.
\begin{enumerate}[label=\textsf{\Alph*.}]
\item \emph{Finite algebra.}  A statement about a finite product, a finite
sum, a coefficient, or a polynomial is translated literally into an
arithmetic identity or inequality between coded integers and rationals.
\item \emph{Effective analysis.}  A real or complex quantity is a fixed
primitive recursive approximation procedure with a fixed primitive recursive
Cauchy modulus.  An equality \(x=y\) means
\[
   \forall p\;\norm{x(p+s)-y(p+s)}\le 2^{-p}
\]
for an explicitly supplied precision shift \(s\).  Uniform statements carry
explicit compact-domain bounds and moduli.
\item \emph{Certified continuation.}  A branch or contour formula is read
relative to a finite code for a path, a branch of logarithm, positive distance
margins from singularities, and quantitative bounds on every local chart.
\item \emph{Global classical claim.}  A statement whose ordinary proof uses
unbounded path existence, compactness without data, an abstract Riemann
surface, or a global classification theorem is not silently reinterpreted as
an iPRA theorem.  We record the precise finite certificate that would suffice,
or mark the claim as outside the present package.
\end{enumerate}
\end{definition}

This distinction makes the strongest part of the theory unexpectedly cheap.
The hypergeometric function is controlled by one rational recurrence, and a
large fraction of its identities are consequences of that recurrence plus
uniqueness of a normalized formal solution of a second-order equation.

\subsection{Principal theorem}

Let a \emph{guarded rational parameter package} consist of rational codes
\(a,b,c,z\), together with proofs that every denominator used below is
nonzero.  A complex-rational package consists of pairs of rational codes and
uses the norm \(\norm{x+iy}=|x|+|y|\), whose comparisons are decidable.
For fixed packages, all operations in the following theorem are primitive
recursive.

\begin{theorem}[Constructive hypergeometric package]\label{thm:main-package}
The following families admit proofs in \(\ipra\), in the senses of
\cref{def:levels}.

\begin{enumerate}[label=(\roman*)]
\item[\StatusA--1.] The rising factorial, hypergeometric coefficient, finite
partial sum, generalized \({}_pF_q\) coefficient, and all their recurrence
relations are primitive recursive after totalized guarded division.

\item[\StatusA--2.] Symmetry in \(a,b\), termination for a nonpositive
integral numerator parameter, the explicit terminating polynomial,
coefficientwise differentiation of every finite order, the hypergeometric
ODE, and uniqueness of its normalized formal solution are provable by
quantifier-free induction.

\item[\StatusA--3.] The pullback--gauge calculations underlying Euler's
and Pfaff's transformations, the standard quadratic transformations, the six
Frobenius expressions, Kummer's 24 local forms, and all contiguous identities
on the source pages reduce to finite polynomial/rational certificates.
Equality of analytically continued branches is read at level \StatusB{} or
\StatusC{} with the displayed branch and path data.  Resonant logarithmic
second solutions are covered to every externally fixed truncation order; a
uniform order requires a jointly primitive recursive coefficient table.

\item[\StatusA--4.] Chu--Vandermonde, the terminating Jacobi and Legendre
representations, the generalized Chebyshev identity as a formal identity, the
binomial, logarithm, arcsine, incomplete-beta coefficient identity, and the
cubic-radical special case are provable by finite algebra and formal-series
uniqueness.

\item[\StatusB--1.] For each rational \(0\le r<1\), the series has a primitive
recursive geometric-tail modulus uniformly on \(\norm z\le r\), provided
quantitative lower bounds are supplied for denominator factors.  The sum is a
primitive recursively named complex number.  Termwise differentiation of any
fixed order, and of variable order under jointly primitive recursive data,
has an explicit modulus.

\item[\StatusB--2.] Euler's beta integral, its \([0,\infty)\) variant,
complete elliptic-integral specializations, and the incomplete-beta formula
hold for effective parameters with positive endpoint-exponent margins and a
branch/apartness certificate.  The proof is a finite quadrature argument plus
a geometric head--tail estimate.

\item[\StatusB--3.] Gauss's summation theorem in its Euler chamber, Kummer's
\(z=-1\) theorem, Gauss's second theorem and Bailey's \(z=1/2\) theorem follow
with explicit error schedules from Euler's integral, beta identities, and
quadratic transformations.  The terminating Vandermonde case has no analytic
hypotheses.

\item[\StatusB--4.] The coefficientwise confluent limit
\[
 {}_1F_1(a;c;z)=\lim_{h\to0}\F{a}{h^{-1}}{c}{hz}
\]
has an explicit compact-disk modulus after splitting into a finite head and a
uniform geometric tail.

\item[\StatusC--1.] The connection formula at \(z=1\), analytic continuation
along a rational polygonal path, Barnes's integral, and Gauss's continued
fraction hold relative to explicit nonresonance, branch, contour-separation,
tail, and denominator-apartness certificates described below.

\item[\StatusD--1.] The global Schwarz-triangle uniformization, the complete
finite-monodromy classification, the modular-lambda inversion statement, and
the modular proofs of all isolated MathWorld algebraic values are not claimed
as completed iPRA reductions here.  We prove the local algebraic differential
identities and give a finite certificate interface for the remaining global
input.
\end{enumerate}
\end{theorem}

\begin{proof}[Organization of the proof]
The finite core is developed in \cref{sec:arithmetic,sec:series,sec:ode}.
Effective convergence and differentiation are proved in
\cref{sec:convergence}.  Elementary and polynomial specializations are in
\cref{sec:special}; transformations and Kummer forms in
\cref{sec:transformations}; integrals and summations in
\cref{sec:integrals-sums}; contiguous relations and continued fractions in
\cref{sec:contiguous}; and continuation, Barnes contours, and the precise
boundary of the package in \cref{sec:continuation}.  Every analytic argument
is accompanied by a primitive recursive choice of truncation, mesh, or
precision.  The Lean compilation plan is in \cref{sec:formalization}.
\end{proof}

\subsection{Why theorem schemas are the right object}

A primitive recursive function may be represented in \(\ipra\) by a function
symbol together with defining equations.  It is important that we do not
quantify over codes for all primitive recursive programs and then assume a
primitive recursive universal evaluator: diagonalization rules that out.
Thus a theorem about a named approximation procedure \(A(n,p)\) is a theorem
schema, one first-order theorem for each fixed primitive recursive definition.
Numerical parameters of that definition may still vary uniformly.

This is exactly the format used by the surrounding ProveIt work.  Its ordinary
object theory has symbols for primitive recursive functions, their defining
equations, successor axioms, and quantifier-free induction.  Its checked
\code{Derivations} library packages theorems as sequents
\[
   ([]:\mathrm{List}(\mathrm{Fm}\;\mathrm{PTm}\;0))
      \;\vdash_{\mathrm{PRAAx}}\;\varphi,
\]
and transports them to the native scalar calculus through a bi-interpretation.
The present report supplies the mathematical derivations and the target
formulas; it does not claim that the full analytic development has already
been compiled into those sequents.

\subsection{Constructive content versus classical notation}

The notation \(\F{a}{b}{c}{z}\) is used at three levels:
\begin{enumerate}
\item the formal series \(\sum_{n\ge0}A_nz^n\), where equality means equality
of every coefficient;
\item the effective sum on a closed disk \(\norm z\le r<1\), represented by
partial sums and a tail modulus;
\item a continued branch obtained from a finite chain of local disks.
\end{enumerate}
The surrounding text always identifies the level.  A formal identity can be
stronger foundationally than an analytic one: it is valid over a localized
polynomial ring with no appeal to limits.  Conversely, a classical formula
that compares branches at distant points needs data that are invisible in the
printed equation.

\begin{remark}[No hidden excluded middle]
All comparisons used to choose a truncation or mesh are comparisons of coded
integers or rationals and hence decidable.  Bounded minimization over such a
test is primitive recursive.  We never infer a modulus from a merely
qualitative convergence assertion, and we never choose a zero or a path by an
unbounded existence principle.
\end{remark}

\subsection{A compact coverage map}

The source pages are covered in the following order.
\begin{center}
\small
\begin{tabularx}{\textwidth}{@{}p{.27\textwidth}p{.12\textwidth}X@{}}
\toprule
Source topic & Status & Constructive replacement \\
\midrule
Series, Pochhammer, termination & \StatusA & finite products, recurrence, finite polynomial \\
Differentiation & \StatusA/\StatusB & coefficient identity; analytic transfer by tail moduli \\
Elementary special cases & \StatusA/\StatusB & formal identities plus effective named functions \\
Hypergeometric ODE & \StatusA & coefficient recurrence and formal uniqueness \\
Six/24 local solutions & \StatusA/\StatusC & algebraic pullback/gauge certificates; branch data for analytic equality \\
Euler integral & \StatusB & rational quadrature with endpoint margins \\
Barnes integral & \StatusC & finite residue rectangle plus explicit vertical/horizontal tail bounds \\
Contiguous relations & \StatusA & coefficientwise rational identities \\
Continued fraction & \StatusA/\StatusC & exact finite convergents; convergence under apartness/contraction data \\
Euler--Pfaff/quadratic transforms & \StatusA/\StatusB & transformed ODE and normalized uniqueness \\
Gauss/Kummer/Bailey sums & \StatusA/\StatusB & finite Vandermonde or beta-integral evaluation \\
Schwarz maps and monodromy & \StatusC/\StatusD & local matrix/path certificates; no global classification claim \\
Modular-lambda and isolated modular values & \StatusD & exact displayed values audited numerically; modular proof not reconstructed \\
\bottomrule
\end{tabularx}
\end{center}

\section{The arithmetic substrate}\label{sec:arithmetic}

\subsection{Integer, rational, and complex-rational codes}

An integer is represented by a pair \((u,v)\in\N^2\), denoting \(u-v\).
A rational is a pair \((m,d)\), where \(m\) is an integer code and \(d+1\)
is the positive denominator.  Equality and order are tested by cross
multiplication.  Normalization by the bounded Euclidean algorithm is useful
for canonical input codes but is not mathematically necessary.  A complex
rational is a pair of rational codes.  We use
\[
  \norm{x+iy}=|x|+|y|
\]
for error estimates.  It satisfies
\(\norm{uv}\le\norm u\norm v\), and every comparison is rational.
The Euclidean norm is used only in familiar displayed formulas; replacing it
by the one-norm changes constants but not the constructions.

Division is totalized.  Define
\[
 \operatorname{qdiv}(x,y)=
 \begin{cases}x/y,&y\ne0,\\0,&y=0.
 \end{cases}
\]
Every theorem that uses the first branch carries a proof of \(y\ne0\).  This
allows all displayed functions to be primitive recursive without turning a
partial operation into an implicit logical principle.

\subsection{Finite sequences and bounded folds}

Pairing and projections code finite tuples.  Bounded sums, products, maxima,
and Boolean folds are obtained by primitive recursion:
\begin{align*}
 \operatorname{sum}(0,f)&=0,&
 \operatorname{sum}(n+1,f)&=\operatorname{sum}(n,f)+f(n),\\
 \operatorname{prod}(0,f)&=1,&
 \operatorname{prod}(n+1,f)&=\operatorname{prod}(n,f)f(n).
\end{align*}
A bounded universal assertion is represented by a Boolean fold
\[
  \operatorname{all}(N,P)=\prod_{n<N}\chi_P(n).
\]
It is equal to \(1\) precisely when every decidable test \(P(n)\) below \(N\)
holds.  This is the standard way to keep induction invariants literally
quantifier-free.  For example, instead of carrying
\(\forall k<n\,D_k\ne0\), the object-theory proof carries
\(\operatorname{all}(n,k\mapsto[D_k\ne0])=1\).

\subsection{Names of effective numbers}

\begin{definition}[Dyadic Cauchy name]\label{def:name}
A complex number name is a primitive recursive function
\(X:\N\to\Q(i)\) satisfying
\[
  \norm{X(p+1)-X(p)}\le2^{-p}
  \qquad(p\in\N).
\]
A stronger and often more convenient contract is
\[
  \norm{X(q)-X(p)}\le2^{-p+1}
  \qquad(q\ge p).
\]
Two names are extensionally equal when a primitive recursive shift \(s\)
satisfies
\[
  \norm{X(p+s)-Y(p+s)}\le 2^{-p}
  \quad\text{for every }p.
\]
\end{definition}

The shift absorbs a finite number of error sources.  Given a rational
constant \(C\ge0\), define by bounded search
\[
 \sigma(C)=\min\{s\le \lceil C\rceil+1:2^s\ge\max(1,C)\}.
\]
Existence follows from \(2^s\ge s+1\), proved by quantifier-free induction.
Then \(C2^{-p-\sigma(C)}\le2^{-p}\).  Every precision allocation below is a
composition of such shifts, bounded searches, and elementary arithmetic, and
is therefore primitive recursive.

\begin{definition}[Named function on a rational domain]\label{def:named-function}
A named function on a coded rational domain \(D\) consists of a primitive
recursive evaluator
\[
  F(x,p)\in\Q(i).
\]
A uniform value bound \(M\), and a modulus \(\omega\) such that
\[
  x,y\in D,\quad \norm{x-y}\le2^{-\omega(p)}
   \Longrightarrow \norm{F(x,p+s)-F(y,p+s)}\le2^{-p}
\]
for a displayed shift \(s\).  Derivative, integral, and branch contracts add
further finite inequalities.  The evaluator is fixed; the theory does not
quantify over all evaluators.
\end{definition}

\subsection{Apartness and branch data}

Constructive division requires more than the negation of equality.  To name
\(1/X\), supply a positive rational margin \(\delta\) and a precision
\(p_0\) such that
\[
  \norm{X(p_0)}\ge\delta,
  \qquad
  \norm{X(p)-X(p_0)}\le\delta/2\quad(p\ge p_0).
\]
Then \(\norm{X}\ge\delta/2\), and rational interval arithmetic gives a name
for the reciprocal.  We call this an \emph{apartness certificate}.

A branch of \(z^\alpha\) on a rational disk consists of a named logarithm
\(L(z)\), a disk center and radius, a proof that the disk avoids the branch
cut and zero by a rational margin, and the local identities
\(L'(z)=1/z\), \(L(z_0)=\ell_0\).  Then
\(z^\alpha=\exp(\alpha L(z))\).  The exponential is represented by its
Taylor series with a factorial-tail modulus.  No global choice of argument is
assumed.

\subsection{Formal power series as finite observations}

A full infinite sequence is not an object over which first-order iPRA
quantifies.  A fixed primitive recursive coefficient function
\(A:\N\to\Q(i)\), however, is legitimate.  A formal-series identity
\[
  \sum_{n\ge0}A_nz^n=\sum_{n\ge0}B_nz^n
\]
means the theorem schema \(A_n=B_n\) for every \(n\).  Product and
composition identities are expressed through finite Cauchy sums.  For
example, the coefficient of a product is
\[
  (A*B)_n=\sum_{k=0}^{n}A_kB_{n-k}.
\]
Every coefficient uses only a bounded sum.

\begin{definition}[Truncated series algebra]\label{def:trunc-series}
For \(N\in\N\), let
\(R_N=R[z]/(z^N)\).  Represent an element by the vector
\((a_0,\ldots,a_{N-1})\).  Addition, Cauchy multiplication, formal
differentiation, and substitution by a series with zero constant term are
primitive recursive operations on finite vectors.  Equality in \(R_N\) is a
finite conjunction.
\end{definition}

A uniform formal identity is proved by an induction whose invariant is the
correctness of all coefficients below the current bound.  Alternatively, for
a fixed target coefficient \(n\), one works entirely in \(R_{n+1}\).  This
``finite observation'' interpretation is especially useful for the Kummer
forms and algebraic substitutions.

\subsection{Isolated algebraic names}

Radicals are handled without postulating a total root operation on arbitrary
reals.

\begin{definition}[Isolated algebraic number]\label{def:alg-name}
An isolated real algebraic name consists of an integer polynomial \(P\),
rational endpoints \(u<v\), a sign-change or monotonicity certificate proving
that \(P\) has exactly one zero in \([u,v]\), and a primitive recursive
bisection procedure.  A complex algebraic name uses two real intervals plus a
finite uniqueness certificate, for example a rectangle on which a resultant
or Newton test is separating.
\end{definition}

Bisection has a primitive recursive convergence rate because the interval
width after \(m\) steps is \((v-u)2^{-m}\).  Square and cube roots occurring
in the source pages are interpreted by positive isolated roots.  Polynomial
identities between such names are proved by interval refinement together with
an exact polynomial relation, not by an unguarded appeal to classical root
choice.

\subsection{Integration by finite quadrature}

For a named uniformly continuous function \(f:[0,1]\to\C\), define the
left Riemann sum
\[
  R_N(f)=\frac1N\sum_{j<N}f(j/N).
\]
If \(\norm{f(x)-f(y)}\le\varepsilon\) whenever
\(|x-y|\le1/N\), then
\[
  \norm{R_N(f)-R_{2N}(f)}\le\varepsilon.
\]
The proof pairs the two fine cells in every coarse cell and uses the supplied
modulus.  Thus a primitive recursive mesh schedule produces a Cauchy name for
the integral.  Endpoint singularities are treated by cutting off
\([0,\eta]\) and \([1-\eta,1]\), bounding those tails by explicit powers,
and applying ordinary quadrature on the central interval.  This is the
integration engine used in \cref{sec:integrals-sums}.

\begin{lemma}[Finite substitution principle]\label{lem:finite-substitution}
Suppose \(f,g\) are named functions on rational intervals, \(g\) is
monotone with named derivative and explicit endpoint values, and the
fundamental theorem estimate is supplied with a modulus.  Then the usual
change-of-variables formula is an iPRA theorem schema.
\end{lemma}

\begin{proof}
Apply the identity to the polygonal/Riemann sums.  On each cell, replace the
increment \(g(x_{j+1})-g(x_j)\) by
\(g'(x_j)(x_{j+1}-x_j)\).  The derivative-remainder contract bounds the
sum of errors by a primitive recursive multiple of the mesh.  The two finite
sums telescope at the endpoints.  Choosing the mesh by bounded search makes
the total error at most \(2^{-p}\).  Every invariant is an equality or
inequality of finite rational sums.
\end{proof}

\section{The finite hypergeometric algebra}\label{sec:series}

\subsection{Rising factorials and coefficients}

\begin{definition}[Rising factorial]\label{def:pochhammer}
For an element \(x\) of a commutative ring, define
\[
  \poch{x}{0}=1,
  \qquad
  \poch{x}{n+1}=\poch{x}{n}(x+n).
\]
For coded rationals or complex rationals this is a primitive recursive finite
product.
\end{definition}

\begin{lemma}[Elementary Pochhammer identities]\label{lem:pochhammer}
Whenever the displayed operations are defined,
\begin{align}
 \poch{x}{m+n}&=\poch{x}{m}\poch{x+m}{n},\label{eq:poch-split}\\
 \poch{x}{n+1}&=x\poch{x+1}{n},\label{eq:poch-shift}\\
 \poch{-m}{n}&=(-1)^n\frac{m!}{(m-n)!}\quad(n\le m),\label{eq:poch-neg}\\
 \poch{-m}{n}&=0\quad(n>m).\label{eq:poch-zero}
\end{align}
All four are provable in iPRA by quantifier-free induction.
\end{lemma}

\begin{proof}
For \eqref{eq:poch-split}, induct on \(n\).  The successor step multiplies
both sides by \(x+m+n\).  Formula \eqref{eq:poch-shift} is the case in which
the first factor is separated.  For \eqref{eq:poch-neg}, write each factor
\(-m+j=-(m-j)\) and induct while maintaining \(n\le m\); cross-multiplication
avoids division.  At \(n=m+1\), the factor \(-m+m\) is zero, and every later
product remains zero.  These are polynomial identities in integer codes.
\end{proof}

\begin{definition}[Gauss coefficient]\label{def:gauss-coeff}
Assume \(\poch{c}{n}\ne0\).  Define
\[
 A_n(a,b;c)=\frac{\poch{a}{n}\poch{b}{n}}
                 {\poch{c}{n}\,n!}.
\]
The guarded total definition uses \(\operatorname{qdiv}\).  The partial sum is
\[
 H_N(a,b;c;z)=\sum_{n<N}A_n(a,b;c)z^n.
\]
\end{definition}

\begin{lemma}[Coefficient ratio]\label{lem:ratio}
Under the denominator guards,
\begin{equation}\label{eq:ratio}
 (n+1)(n+c)A_{n+1}=(n+a)(n+b)A_n.
\end{equation}
If \((n+1)(n+c)\) is apart from zero, this is equivalently
\[
  \frac{A_{n+1}}{A_n}=\frac{(n+a)(n+b)}{(n+1)(n+c)}
\]
when \(A_n\ne0\).
\end{lemma}

\begin{proof}
Expand the four successor recurrences for the rising factorials and factorial,
then cancel only after multiplying through by the denominators.  The
cross-multiplied identity \eqref{eq:ratio} is valid even when \(A_n=0\).
\end{proof}

\begin{definition}[Formal Gauss series]\label{def:formal-gauss}
The formal object
\[
   \F{a}{b}{c}{z}:=\sum_{n\ge0}A_n(a,b;c)z^n
\]
means the primitive recursive coefficient table \(n\mapsto A_n\), together
with the guards \(\poch c n\ne0\) needed at the observed indices.
\end{definition}

\begin{proposition}[Symmetry and normalization]\label{prop:symmetry}
Formally,
\[
 \F{a}{b}{c}{z}=\F{b}{a}{c}{z},
 \qquad
 \F{a}{b}{c}{0}=1.
\]
\end{proposition}

\begin{proof}
The coefficient formula is symmetric in \(a,b\), and the constant coefficient
is \(1\).  No limit is involved.
\end{proof}

\subsection{Termination and the polynomial formula}

\begin{theorem}[Terminating series]\label{thm:termination}
For \(m\in\N\),
\begin{equation}\label{eq:terminating}
 \F{-m}{b}{c}{z}
   =\sum_{n=0}^{m}(-1)^n\binom mn\frac{\poch b n}{\poch c n}z^n
\end{equation}
as a polynomial identity, provided the finitely many denominator factors are
nonzero.
\end{theorem}

\begin{proof}
By \cref{lem:pochhammer},
\[
 \frac{\poch{-m}{n}}{n!}
 =(-1)^n\frac{m!}{n!(m-n)!}=(-1)^n\binom mn
 \quad(n\le m),
\]
and every coefficient with \(n>m\) vanishes.  Cross-multiplying by the finite
product \(\prod_{j<m}(c+j)\) turns the whole statement into an integer
polynomial identity.  The proof uses bounded products and induction only.
\end{proof}

\begin{corollary}[Binomial specialization]\label{cor:binomial-termination}
For \(m\in\N\),
\[
 \F{-m}{b}{b}{z}=(1-z)^m.
\]
In particular, \((1+z)^m=\F{-m}{b}{b}{-z}\).
\end{corollary}

\begin{proof}
The factors \(\poch b n\) cancel under the supplied guards, and
\eqref{eq:terminating} becomes the finite binomial theorem.  The binomial
theorem itself is a quantifier-free induction on \(m\).
\end{proof}


\subsection{Regularization at a singular denominator parameter}

The ordinary series is singular when \(c\) approaches a nonpositive integer,
but multiplication by the reciprocal gamma function removes the singularity.

\begin{theorem}[Regularized limit at \(c=-m\)]\label{thm:regularized-c}
Let \(m\in\N\).  Coefficientwise, and uniformly on every strict disk after
supplying the effective reciprocal-gamma bounds,
\begin{equation}\label{eq:regularized-c}
 \lim_{c\to-m}\frac{\F a b c z}{\Gamma(c)}
 =\frac{\poch a{m+1}\poch b{m+1}}{(m+1)!}z^{m+1}
  \F{a+m+1}{b+m+1}{m+2}{z}.
\end{equation}
\end{theorem}

\begin{proof}
Use the exact gamma recurrence identity
\[
 \frac1{\Gamma(c)\poch c n}=\frac1{\Gamma(c+n)}.
\]
Thus the coefficient of \(z^n\) on the left is
\[
 \frac{\poch a n\poch b n}{n!\,\Gamma(c+n)}.
\]
The reciprocal gamma function has a zero at every nonpositive integer.  At
\(c=-m\), the coefficients with \(n\le m\) therefore tend to zero.  Write
\(n=m+1+k\) for the remaining coefficients.  Since
\(\Gamma(k+1)=k!\), their limits are
\[
 \frac{\poch a{m+1+k}\poch b{m+1+k}}
      {(m+1+k)!\,k!}.
\]
Splitting both rising factorials at \(m+1\) and using
\((m+2)_k=(m+1+k)!/(m+1)!\) gives exactly the coefficient on the right of
\eqref{eq:regularized-c}.  This is a finite identity for each coefficient.
For uniform convergence on \(\norm z\le r<1\), choose a small rational
parameter disk around \(-m\), a common ratio cutoff for its tail, and then a
finite head on which the reciprocal-gamma name converges effectively.  The
gamma packet supplies the simple-zero and local modulus data.
\end{proof}

\subsection{Generalized hypergeometric coefficients}

\begin{definition}[Generalized coefficient]\label{def:pfq}
For finite parameter vectors
\(\boldsymbol a=(a_1,\ldots,a_p)\) and
\(\boldsymbol b=(b_1,\ldots,b_q)\), put
\[
 T_n(\boldsymbol a;\boldsymbol b;z)
  =\frac{\prod_{i<p}\poch{a_i}{n}}
         {n!\prod_{j<q}\poch{b_j}{n}}z^n.
\]
The associated formal series is \({}_pF_q\).  All vector products are bounded
folds over the externally fixed or coded finite lengths \(p,q\).
\end{definition}

\begin{proposition}[Hypergeometric term ratio]\label{prop:pfq-ratio}
The successive terms satisfy the cross-multiplied identity
\[
 (n+1)\prod_{j<q}(n+b_j)T_{n+1}
 =z\prod_{i<p}(n+a_i)T_n.
\]
Thus the ratio of successive nonzero terms is a rational function of \(n\),
as asserted on the MathWorld page.
\end{proposition}

\begin{proof}
Apply the successor equation for every Pochhammer factor.  Bounded-product
congruence and associativity suffice.
\end{proof}

\subsection{Formal differentiation}

For a formal series \(F(z)=\sum_{n\ge0}f_nz^n\), define
\(DF=\sum_{n\ge0}(n+1)f_{n+1}z^n\).  This is a primitive recursive operation
on coefficient tables.

\begin{theorem}[Derivative formulas]\label{thm:derivatives}
For every \(r\in\N\) for which the denominator guards hold,
\begin{equation}\label{eq:nth-derivative}
 D^r\F{a}{b}{c}{z}
 =\frac{\poch a r\poch b r}{\poch c r}
   \F{a+r}{b+r}{c+r}{z}.
\end{equation}
The theorem is uniform in \(r\) when the shifted coefficient table and guards
are jointly primitive recursive.
\end{theorem}

\begin{proof}
For \(r=1\), the coefficient of \(z^n\) on the left is
\[
 (n+1)\frac{\poch a{n+1}\poch b{n+1}}
                   {\poch c{n+1}(n+1)!}
 =\frac{ab}{c}\frac{\poch{a+1}n\poch{b+1}n}
                         {\poch{c+1}n n!},
\]
using \eqref{eq:poch-shift}.  This is an exact rational identity.  Induction
on \(r\), together with
\(\poch a{r+1}=\poch a r(a+r)\), gives
\eqref{eq:nth-derivative}.  The induction invariant is a single
cross-multiplied coefficient equality.
\end{proof}

\subsection{The coefficient recurrence as a characterization}

\begin{theorem}[Formal uniqueness]\label{thm:formal-uniqueness}
Let \((B_n)\) be a coefficient table satisfying
\[
 B_0=1,
 \qquad
 (n+1)(n+c)B_{n+1}=(n+a)(n+b)B_n,
\]
and suppose each \((n+1)(n+c)\) up to the observed index is apart from zero.
Then \(B_n=A_n(a,b;c)\) for every \(n\).
\end{theorem}

\begin{proof}
Induct on \(n\).  The base is \(1\).  In the successor step, both \(A_{n+1}\)
and \(B_{n+1}\) solve the same linear equation with an apart coefficient.
Cross-multiply and use cancellation in the coded rational field.  In the
object theory one may avoid field cancellation by carrying numerator and
denominator representatives and the nonzero product as a side assumption.
\end{proof}

\begin{remark}[Why uniqueness is powerful]
To prove a transformation, it often suffices to show that a transformed
series has constant term \(1\) and satisfies the same recurrence.  This turns
an identity between infinite objects into a finite calculation with the
chain rule, followed by \cref{thm:formal-uniqueness}.
\end{remark}

\subsection{Finite verification principle}

\begin{proposition}[Certificate checker for identities]\label{prop:certificate-checker}
Fix rational parameter maps
\((a,b,c,z)\mapsto(a',b',c',\phi(z))\) and a prefactor \(P(z)\) whose formal
coefficients are primitive recursive.  To prove
\[
 \F a b c z=P(z)\F{a'}{b'}{c'}{\phi(z)}
\]
it is enough to provide either:
\begin{enumerate}
\item a coefficient identity for the \(n\)-th coefficient, expressed as a
finite sum and proved by induction in \(n\); or
\item an algebraic certificate that the right side satisfies the same
second-order differential equation and the same first two initial
coefficients as the left side.
\end{enumerate}
Both certificate formats have primitive recursive checkers.
\end{proposition}

\begin{proof}
In the first format, the checker expands products and substitutions in
\(R_{n+1}\).  In the second, the differential expression has coefficients
obtained by a finite number of Cauchy products and derivatives; equality to
zero at each coefficient is decidable.  The ODE recurrence then determines
all higher coefficients from the first two, or from the constant coefficient
for the normalized regular solution.  The checker never manipulates an
actual completed infinite series.
\end{proof}

\section{Effective convergence and differentiation}\label{sec:convergence}

\subsection{A primitive recursive ratio cutoff}

Fix complex-rational parameters \(a,b,c\) and a rational radius
\(0\le r<1\).  Let rational numbers \(A,B,C\) satisfy
\[
 \norm a\le A,\qquad \norm b\le B,\qquad \norm c\le C.
\]
For \(n>C\), the reverse triangle inequality gives
\(\norm{n+c}\ge n-C\).  Hence, for \(\norm z\le r\),
\begin{equation}\label{eq:term-ratio-bound}
 \norm{\frac{A_{n+1}z^{n+1}}{A_nz^n}}
 \le r\frac{(n+A)(n+B)}{(n+1)(n-C)}
 \le r\frac{(n+A)(n+B)}{n(n-C)}.
\end{equation}
This estimate is used only after denominator apartness has been verified.

\begin{lemma}[Explicit ratio cutoff]\label{lem:ratio-cutoff}
If \(0<r<1\), put
\[
 q=\frac{1+r}{2},\qquad
 \varepsilon=\frac{q-r}{r}=\frac{1-r}{2r},
\]
\[
 D=(1+\varepsilon)C+A+B+AB,
 \qquad
 N_0=1+\left\lceil\max\left(C,\frac{2D}{\varepsilon}\right)\right\rceil.
\]
Then for every \(n\ge N_0\), the right side of
\eqref{eq:term-ratio-bound} is at most \(q\).  For \(r=0\), every nonconstant
term vanishes.
\end{lemma}

\begin{proof}
It is enough to prove
\[
 (n+A)(n+B)\le(1+\varepsilon)n(n-C).
\]
After expansion this is
\[
 \varepsilon n^2\ge
  ((1+\varepsilon)C+A+B)n+AB.
\]
For \(n\ge1\), the right side is at most \(Dn\).  The definition of
\(N_0\) gives \(\varepsilon n\ge2D\), hence the desired inequality with
room to spare.  Every comparison is rational and every ceiling is obtained by
bounded integer search.
\end{proof}

\begin{theorem}[Uniform convergence on a strict subdisk]\label{thm:disk-convergence}
Assume the denominator factors \(c+j\) carry apartness certificates.  For
every rational \(r<1\), the partial sums
\(H_N(a,b;c;z)\) have a primitive recursive Cauchy modulus uniformly on
\(\norm z\le r\).  More precisely, for \(N\ge N_0\) and \(m\ge0\),
\begin{equation}\label{eq:geometric-tail}
 \norm{\sum_{n\ge N+m}A_nz^n}
 \le \frac{\norm{A_N}r^N}{1-q}q^m.
\end{equation}
A primitive recursive bounded search finds the least \(m\) making the right
side at most \(2^{-p}\).
\end{theorem}

\begin{proof}
By \cref{lem:ratio-cutoff}, every term after \(N_0\) is at most \(q\) times
the preceding term.  Induction gives
\(\norm{A_{N+j}z^{N+j}}\le\norm{A_N}r^Nq^j\).
The finite geometric identity
\[
 (1-q)\sum_{j<m}q^j=1-q^m
\]
is an exact polynomial identity, so the tail bound follows by passing to the
Cauchy estimate.  No classical ratio test is invoked: the ratio cutoff and
the geometric remainder are explicit.

For termination of the search, write \(d=1-q>0\).  The finite binomial
inequality
\((1-d)^{-m}\ge1+md\) gives
\(q^m\le(1+md)^{-1}\).  Taking
\(m\ge2^{p+s}/d\), where \(s\) absorbs
\(\norm{A_N}r^N/(1-q)\), certainly succeeds.  This is a primitive recursive,
though intentionally crude, bound.
\end{proof}

\begin{corollary}[Named analytic function on a strict subdisk]
For each fixed guarded parameter package and rational \(r<1\),
\(z\mapsto\F a b c z\) is a named uniformly continuous function on
\(\norm z\le r\).  An evaluator computes a sufficiently long rational
partial sum; a modulus follows by splitting
\[
 F(z)-F(w)=\sum_{n<N}A_n(z^n-w^n)+\Tail_N(z)-\Tail_N(w).
\]
\end{corollary}

\begin{proof}
Choose \(N\) so that both tails are at most \(2^{-p-2}\).  The finite head is
uniformly continuous because
\[
 z^n-w^n=(z-w)\sum_{j<n}z^jw^{n-1-j}.
\]
On \(\norm z,\norm w\le r\), its norm is at most
\(n r^{n-1}\norm{z-w}\).  The finite sum of these rational constants gives an
explicit Lipschitz bound and hence a dyadic modulus.
\end{proof}

\subsection{Analytic differentiation}

The formal identity of \cref{thm:derivatives} becomes an analytic identity
once both sides have compatible moduli.

\begin{theorem}[Effective termwise differentiation]\label{thm:analytic-derivative}
Let \(0\le r<\rho<1\) be rational.  On \(\norm z\le r\), the difference
quotient of the named sum converges effectively to the named derivative
series, and
\[
 \frac{d^s}{dz^s}\F a b c z
 =\frac{\poch a s\poch b s}{\poch c s}
   \F{a+s}{b+s}{c+s}{z}.
\]
after the shifted denominator guards are supplied.
\end{theorem}

\begin{proof}
Choose the series on the larger disk \(\norm z\le\rho\).  Its tail has a
primitive recursive modulus by \cref{thm:disk-convergence}.  For a monomial,
\[
 \frac{(z+h)^n-z^n}{h}-nz^{n-1}
 =\sum_{j=2}^{n}\binom nj z^{n-j}h^{j-1}.
\]
On \(\norm z\le r\) and \(\norm h\le(\rho-r)/2\), the norm is bounded by an
explicit finite multiple of \(\norm h\rho^{n-2}\).  Split the derivative
series into a finite head and tail.  Choose the tail cutoff first, then choose
\(h\) small enough for the finite head.  All choices are bounded searches over
rational inequalities.  The coefficient identity from
\cref{thm:derivatives} identifies the resulting derivative name.  Iterating
this argument gives fixed order \(s\); joint primitive recursive data give a
uniform variable-order schema.
\end{proof}

\subsection{Convergence on the unit circle}

The source pages state convergence at \(|z|=1\) when
\(\RePart(c-a-b)>0\).  A constructive version needs a positive margin, not
merely the negation of a nonpositive real part.

\begin{contract}[Boundary convergence data]\label{con:boundary}
Supply a positive rational \(\delta\) with
\[
 \RePart(c-a-b)\ge\delta>0,
\]
a denominator-apartness table, and an effective gamma-ratio estimate
\begin{equation}\label{eq:gamma-ratio-contract}
 \norm{A_n(a,b;c)}\le K(n+1)^{-1-\delta}
 \qquad(n\ge N_1)
\end{equation}
for rational \(K,N_1\).  The gamma/Stirling package in the repository gives a
route to such an estimate for fixed primitive recursive parameter names on a
compact nonresonant set.
\end{contract}

\begin{theorem}[Effective unit-circle convergence]\label{thm:unit-circle}
Under \cref{con:boundary}, the hypergeometric series converges absolutely and
uniformly on the full coded Euclidean unit disk
\[
   (\RePart z)^2+(\ImPart z)^2\le1,
\]
with a primitive recursive tail modulus.  In particular it converges
uniformly on the rational diamond \(\norm z\le1\).
\end{theorem}

\begin{proof}
Choose an integer \(k\ge1\) with \(1/k\le\delta\).  For \(N\ge N_1\), group
indices into blocks
\([j^k,(j+1)^k)\).  Every term in the block is at most
\(Kj^{-k-1}\), while the block has at most
\(k(j+1)^{k-1}\) terms.  Hence its contribution is at most
\(K'j^{-2}\) for an explicit rational \(K'\).  The elementary telescoping
bound
\[
 \sum_{j\ge J}\frac1{j^2}
 \le\sum_{j\ge J}\frac1{j(j-1)}=\frac1{J-1}
\]
gives a primitive recursive tail cutoff.  This avoids invoking an integral
test or a completed real-number order.  If
\((\RePart z)^2+(\ImPart z)^2\le1\), then the same inequality, raised to the
\(n\)-th power, certifies that each real and imaginary component of \(z^n\)
has absolute value at most one.  Consequently
\(\norm{A_nz^n}\le2\norm{A_n}\); the harmless factor two is absorbed in
\(K'\).  On \(\norm z\le1\), even that extra argument is unnecessary.
\end{proof}

\begin{remark}[Endpoint without a margin]
The qualitative assertion \(\RePart(c-a-b)>0\) is read as the existence of a
positive rational witness \(\delta\).  For a fixed named real with no supplied
witness, iPRA cannot in general decide on which side of zero it lies.  The
margin is mathematical data, not bureaucratic overhead.
\end{remark}

\subsection{Coefficientwise confluent limit}

Let \(h\ne0\) and consider
\[
 G_h(z)=\F a{h^{-1}}c{hz}.
\]
A coefficient contains
\[
 \poch{h^{-1}}n h^n=\prod_{j<n}(1+jh).
\]

\begin{theorem}[Effective confluent limit]\label{thm:confluent}
For every fixed \(n\),
\[
 \prod_{j<n}(1+jh)\longrightarrow1
\]
with the explicit bound
\begin{equation}\label{eq:finite-product-perturb}
 \norm{\prod_{j<n}(1+jh)-1}
 \le \left(1+n\norm h\right)^n-1.
\end{equation}
Consequently the coefficients of \(G_h\) converge to those of
\[
 {}_1F_1(a;c;z)=\sum_{n\ge0}\frac{\poch a n}{\poch c n}\frac{z^n}{n!}.
\]
On each rational disk \(\norm z\le r\), and for \(h\) restricted by a
sufficiently small rational bound, the convergence is uniform with a primitive
recursive modulus.
\end{theorem}

\begin{proof}
Expand the difference of two finite products by telescoping one factor at a
time:
\[
 \prod_{j<n}(1+u_j)-1
 =\sum_{k<n}u_k\prod_{j<k}(1+u_j).
\]
With \(u_j=jh\), a coarse estimate gives
\eqref{eq:finite-product-perturb}.  For uniform convergence, choose a head
cutoff \(N\) using a common ratio majorant for all sufficiently small
\(\norm h\).  On the finite head, choose \(h\) so that each coefficient error
is below the allocated budget.  The two tails use the same geometric bound.
This is the constructive head--tail proof; no interchange of two abstract
limits is needed.
\end{proof}

\section{The differential equation and its local solutions}\label{sec:ode}

\subsection{Coefficient proof of the equation}

Let
\[
 \mathcal L_{a,b,c}y
 =z(1-z)y''+[c-(a+b+1)z]y'-aby.
\]

\begin{theorem}[Hypergeometric differential equation]\label{thm:ode}
The formal Gauss series satisfies
\begin{equation}\label{eq:hg-ode}
 \mathcal L_{a,b,c}\F a b c z=0.
\end{equation}
On every strict subdisk where the series is effectively convergent, the same
identity holds between named analytic functions.
\end{theorem}

\begin{proof}
For \(y=\sum_{n\ge0}A_nz^n\), the coefficient of \(z^n\) in
\(\mathcal L_{a,b,c}y\) is
\[
 (n+1)(n+c)A_{n+1}
 -[n^2+(a+b)n+ab]A_n
 =(n+1)(n+c)A_{n+1}-(n+a)(n+b)A_n.
\]
It vanishes by \cref{lem:ratio}.  This is a single rational identity for
each \(n\).  Analytic equality follows from
\cref{thm:analytic-derivative}: all three named series have explicit common
tail moduli, so coefficientwise cancellation passes to their sums.
\end{proof}

\begin{corollary}[Normalized regular solution is unique]\label{cor:ode-unique}
Suppose \(y=\sum B_nz^n\), \(B_0=1\), and
\(\mathcal L_{a,b,c}y=0\).  If \(n+c\) is apart from zero for every observed
\(n\), then \(y=\F a b c z\) formally.
\end{corollary}

\begin{proof}
The ODE coefficient gives exactly the recurrence in
\cref{thm:formal-uniqueness}.
\end{proof}

\subsection{Frobenius calculation at the three singular points}

The following calculations are algebra in a ring with a formal monomial
\(z^s\) and the derivative rule
\(D(z^s)=sz^{s-1}\).  Analytically, the monomial is interpreted using a
certified logarithm branch.

\begin{proposition}[Solutions at zero]\label{prop:zero-solutions}
The indicial polynomial at \(z=0\) is
\[
 I_0(s)=s(s+c-1).
\]
The two formal Frobenius expressions are
\begin{align}
 u_{0,1}(z)&=\F a b c z,\label{eq:u01}\\
 u_{0,2}(z)&=z^{1-c}
   \F{a+1-c}{b+1-c}{2-c}{z}.
   \label{eq:u02}
\end{align}
Under nonresonance and branch certificates they are named local analytic
solutions.
\end{proposition}

\begin{proof}
Substitute \(y=z^sv\).  The coefficient of the lowest power is
\(s(s+c-1)v(0)\).  For \(s=0\), the equation for \(v\) is the original one.
For \(s=1-c\), direct collection gives
\[
 \mathcal L_{a,b,c}(z^{1-c}v)
 =z^{1-c}\mathcal L_{a+1-c,b+1-c,2-c}v.
\]
The claimed solutions now follow from \cref{thm:ode}.  The displayed gauge
identity is checked after clearing the monomial by a finite polynomial
calculation in \(s\); substituting \(s=1-c\) finishes it.
\end{proof}

\begin{proposition}[Solutions at one]\label{prop:one-solutions}
With \(t=1-z\), the exponents at \(z=1\) are \(0\) and \(c-a-b\), and two
standard expressions are
\begin{align}
 u_{1,1}(z)&=\F a b{a+b+1-c}{1-z},\label{eq:u11}\\
 u_{1,2}(z)&=(1-z)^{c-a-b}
  \F{c-a}{c-b}{c+1-a-b}{1-z}.
  \label{eq:u12}
\end{align}
\end{proposition}

\begin{proof}
Replace \(z\) by \(1-t\) and use \(d/dz=-d/dt\).  After collecting terms,
the regular equation in \(t\) has parameters
\((a,b,a+b+1-c)\).  The second expression is obtained from the second
exponent exactly as in \cref{prop:zero-solutions}.  Every step is an identity
of rational differential operators.
\end{proof}

\begin{proposition}[Solutions at infinity]\label{prop:infty-solutions}
With \(w=z^{-1}\), the exponents at infinity are \(a,b\) in the convention
\(y\sim z^{-s}\).  Two expressions are
\begin{align}
 u_{\infty,a}(z)&=(-z)^{-a}
   \F a{a+1-c}{a+1-b}{z^{-1}},\label{eq:uinf-a}\\
 u_{\infty,b}(z)&=(-z)^{-b}
   \F{b+1-c}{b}{b+1-a}{z^{-1}}.
   \label{eq:uinf-b}
\end{align}
The choice of \((-z)^{-a}\) rather than \(z^{-a}\) is branch bookkeeping;
formally either monomial is a unit multiple of the other.
\end{proposition}

\begin{proof}
Use
\(D_z=-w^2D_w\) and
\(D_z^2=w^4D_w^2+2w^3D_w\), then set
\(y=z^{-s}v(w)\).  The lowest-order equation is
\((s-a)(s-b)=0\).  At \(s=a\) and \(s=b\), collection gives the two displayed
parameter triples.  A finite symbolic checker can verify the operator
identity after multiplying by a sufficient power of \(z\).
\end{proof}

\begin{remark}[Resonance]
If an exponent difference is an integer, the second local solution may
contain a logarithm.  The nonlogarithmic formulas can coincide or become
undefined because a denominator parameter is a nonpositive integer.  The
constructive response is not to pretend that the generic basis persists.
For every fixed truncation order, substitute
\[
 y=y_1\log z+z^s\sum_{n<N}d_nz^n+O(z^{s+N})
\]
into the differential equation.  The coefficients \(d_n\) satisfy a finite
inhomogeneous recurrence involving finite harmonic sums (equivalently,
digamma differences).  This recurrence is primitive recursive.  A full
uniform logarithmic solution is a \StatusB{} theorem once a joint coefficient
bound and convergence modulus are supplied.  The present package gives the
recurrence interface but does not claim a completed all-resonances analytic
library.
\end{remark}

\subsection{Wronskians and local independence}

\begin{lemma}[Wronskian equation]\label{lem:wronskian}
For two solutions of \eqref{eq:hg-ode}, their Wronskian satisfies
\[
 W'(z)=-\frac{c-(a+b+1)z}{z(1-z)}W(z).
\]
Equivalently
\[
 W(z)=Cz^{-c}(1-z)^{c-a-b-1}
\]
on a certified branch domain.
\end{lemma}

\begin{proof}
For an equation \(y''+P y'+Qy=0\), direct expansion gives
\(W'=-PW\).  Here
\[
 P(z)=\frac{c-(a+b+1)z}{z(1-z)}
      =\frac c z+\frac{c-a-b-1}{1-z}.
\]
The displayed solution follows by differentiating the right side.  Thus the
claim can be verified without invoking indefinite integration.  A nonzero
constant \(C\) is witnessed by the leading terms of a nonresonant Frobenius
pair.
\end{proof}

\subsection{The normal or Q-form}

\begin{proposition}[Elimination of the first derivative]\label{prop:q-form}
Set
\[
 y(z)=v(z)u(z),
 \qquad
 v(z)=z^{-c/2}(1-z)^{(c-a-b-1)/2}.
\]
Then \(u\) satisfies
\[
 u''+Q(z)u=0.
\]
After choosing the convention consistent with \(y=vu\), where
\[
 Q(z)=
 \frac{z^2[1-(a-b)^2]
       +z[2c(a+b-1)-4ab]+c(2-c)}
      {4z^2(1-z)^2}.
\]
\end{proposition}

\begin{proof}
Divide the original equation by \(z(1-z)\), obtaining
\(y''+Py'+Ry=0\).  For \(y=vu\), the coefficient of \(u'\) is
\(2v'/v+P\), which vanishes because
\[
 \frac{v'}v=-\frac{c}{2z}-\frac{1+a+b-c}{2(z-1)}=-\frac P2.
\]
The remaining coefficient is \(R-P'/2-P^2/4\).  Clearing the denominator
\(4z^2(1-z)^2\) and expanding yields the displayed numerator.  This is a
polynomial identity suitable for the repository's ring tactic.
\end{proof}

\subsection{Reduction of a three-singularity equation}

The classical slogan that every second-order linear equation with three
regular singular points is hypergeometric must be made data-sensitive.

\begin{theorem}[Certified Riemann-to-Gauss reduction]\label{thm:riemann-reduction}
Let
\[
 y''+P(x)y'+Q(x)y=0
\]
be a rational differential equation.  Supply:
\begin{enumerate}
\item three pairwise apart algebraic points \(x_0,x_1,x_\infty\), with a
certificate that the poles of \(P,Q\) occur only there;
\item regular-singularity bounds
\((x-x_i)P(x)\) and \((x-x_i)^2Q(x)\) at each finite point, and the analogous
bounds at infinity;
\item isolated roots of the three indicial polynomials;
\item the Fuchs relation between the six exponents.
\end{enumerate}
Then a primitive recursive symbolic procedure constructs a fractional-linear
coordinate \(z\) sending the three points to \(0,1,\infty\), a gauge
\(y=z^\alpha(1-z)^\beta u\), and parameters \(a,b,c\) for which \(u\)
satisfies the hypergeometric equation.  The correctness proof is a finite
rational-function identity.
\end{theorem}

\begin{proof}
Take
\[
 z=\frac{(x-x_0)(x_1-x_\infty)}
          {(x-x_\infty)(x_1-x_0)}.
\]
Apartness makes every denominator invertible.  The chain rule transforms the
equation to one with singularities at \(0,1,\infty\).  Choose the gauge
exponents to make one exponent at \(0\) and one at \(1\) equal to zero.  If
the remaining exponent differences are \(1-c\), \(c-a-b\), and \(a-b\),
solve the resulting linear equations for \(a,b,c\).  The Fuchs relation
ensures consistency.  After multiplication by a common polynomial
denominator, equality with \(\mathcal L_{a,b,c}\) is an exact coefficient
identity.  All root choices were supplied as isolated algebraic names, so no
unbounded selection is hidden in the construction.
\end{proof}

\section{Elementary, algebraic, and polynomial special cases}\label{sec:special}

\subsection{The binomial and geometric series}

\begin{theorem}[Generalized binomial identity]\label{thm:binomial}
Formally,
\begin{equation}\label{eq:binomial}
 \F a b b z=\sum_{n\ge0}\frac{\poch a n}{n!}z^n=(1-z)^{-a}.
\end{equation}
For rational or algebraically named \(a\), the right side means the normalized
formal solution of
\((1-z)y'-ay=0\).  On a certified branch disk it is the named analytic power
\(\exp(-a\log(1-z))\).
\end{theorem}

\begin{proof}
The cancellation of \(\poch b n\) is finite and guarded.  If
\(B_n=\poch a n/n!\), then
\((n+1)B_{n+1}=(n+a)B_n\), which is exactly the coefficient recurrence of
\((1-z)y'-ay=0\), with \(B_0=1\).  The named power has the same differential
equation and initial value; effective uniqueness on a strict disk follows by
integrating the difference or by comparing Taylor coefficients and tails.
\end{proof}

\begin{corollary}[Geometric cases]
\begin{align*}
 \F{1}{b}{b}{z}&=\frac1{1-z}, &
 \F{1}{1}{1}{z}&=\frac1{1-z},\\
 \F{1}{2}{1}{z}&=\frac1{(1-z)^2}, &
 \F{1}{2}{2}{z}&=\frac1{1-z}.
\end{align*}
\end{corollary}

\begin{proof}
Apply \cref{thm:binomial}; for \(\F12 1z\), the coefficient is \(n+1\), and
the finite identity
\((1-z)^2\sum_{n<N}(n+1)z^n=1-(N+1)z^N+Nz^{N+1}\) supplies an explicit
remainder.
\end{proof}

\subsection{Logarithm and inverse sine}

\begin{theorem}[Logarithm]\label{thm:log-special}
On \(\norm z<1\), with the branch normalized by \(\log(1)=0\),
\[
 \F11 2{-z}=\frac{\log(1+z)}z,
\]
with the value at \(z=0\) interpreted as \(1\).
\end{theorem}

\begin{proof}
The coefficient is
\[
 \frac{\poch1n^2}{\poch2n n!}=\frac1{n+1}.
\]
Thus
\[
 z\F11 2{-z}=\sum_{m\ge1}(-1)^{m-1}\frac{z^m}{m}.
\]
Formal differentiation gives \((1+z)^{-1}\) and the constant term is zero.
Define the effective logarithm either by this series or by
\(\int_0^z(1+t)^{-1}\,dt\) on a star-shaped certified disk.  The geometric
tail proof in \cref{sec:convergence} makes the equality quantitative.
\end{proof}

\begin{theorem}[Arcsine]\label{thm:arcsin-special}
On a certified disk \(\norm z<1\),
\[
 \F{1/2}{1/2}{3/2}{z^2}=\frac{\arcsin z}{z},
\]
with value \(1\) at zero.
\end{theorem}

\begin{proof}
The finite factorial identities
\[
 \frac{\poch{1/2}n}{n!}=\frac1{4^n}\binom{2n}{n},
 \qquad
 \frac{\poch{1/2}n}{\poch{3/2}n}=\frac1{2n+1}
\]
are proved by induction after clearing powers of two.  Hence the coefficient
of \(z^{2n}\) is
\[
 \frac1{4^n(2n+1)}\binom{2n}{n}.
\]
Multiplying by \(z\) and differentiating gives the binomial series
\((1-z^2)^{-1/2}\), while the value at zero is zero.  This is the defining
effective integral for \(\arcsin z\) on the chosen branch.
\end{proof}

\subsection{The cubic radical identity}

The most striking elementary formula on the Wikipedia page can be proved
without guessing the radical: both sides are the uniquely normalized algebraic
solution of a cubic.

\begin{theorem}[Cubic algebraic hypergeometric function]\label{thm:cubic-radical}
Let
\[
 Y(x)=\F{1/3}{2/3}{3/2}{-27x^2/4}.
\]
As a formal power series,
\begin{equation}\label{eq:cubic-algebraic}
 Y+x^2Y^3=1,
 \qquad Y(0)=1.
\end{equation}
On the real branch with positive square and cube roots,
\begin{equation}\label{eq:cubic-radical}
Y(x)=
\frac{
 \sqrt[3]{\frac{3\sqrt3x+\sqrt{27x^2+4}}2}
 -\sqrt[3]{\frac{2}{3\sqrt3x+\sqrt{27x^2+4}}}
}{\sqrt3x},
\end{equation}
with the removable value \(1\) at \(x=0\).
\end{theorem}

\begin{proof}
Write \(t=x^2\).  Transforming the hypergeometric ODE gives
\begin{equation}\label{eq:cubic-ode}
 t(4+27t)Y''+(6+54t)Y'+6Y=0,
 \qquad Y(0)=1,\quad Y'(0)=-1.
\end{equation}
Now let \(U\) be the unique formal solution of \(U+tU^3=1\) with
\(U(0)=1\).  Existence and uniqueness are coefficientwise: the coefficient of
\(t^n\) in the cubic determines \(u_n\) from earlier coefficients.  Differentiate
the cubic twice, eliminate \(U'\) and \(U''\), and reduce with
\(tU^3=1-U\).  Exact polynomial simplification yields
\eqref{eq:cubic-ode}.  Its first two coefficients agree with those of \(Y\),
so the ODE recurrence gives \(U=Y\).  This entire calculation can be checked
in the finite truncated-series ring of \cref{def:trunc-series}.

For the radical, set
\[
 s=\sqrt{27x^2+4},\qquad
 u^3=\frac{s+3\sqrt3x}{2},\qquad v=u-u^{-1}.
\]
Because \((s+3\sqrt3x)(s-3\sqrt3x)=4\), one has
\(u^{-3}=(s-3\sqrt3x)/2\), and therefore
\[
 v^3+3v=u^3-u^{-3}=3\sqrt3x.
\]
Putting \(R=v/(\sqrt3x)\) gives \(R+x^2R^3=1\).  The positive isolated-root
choices make \(R\to1\) at zero, so formal/algebraic uniqueness gives
\(R=Y\).  Bisection names for the radicals and the strict-disk series modulus
turn the equality into an iPRA theorem schema.
\end{proof}

\subsection{A general algebraic family (Glasser's identity)}

\begin{theorem}[Well-poised algebraic family]\label{thm:glasser}
Let \(k\ge2\),
\[
 f_k=\frac{k}{(k+1)^{1+1/k}}.
\]
The numerator parameter vector
\(\frac1{k+1},\frac2{k+1},\ldots,\frac{k}{k+1}\), and denominator vector
\(\frac2k,\frac3k,\ldots,\frac{k-1}k,\frac{k+1}k\).  Then, as a formal
identity in \(x\),
\begin{equation}\label{eq:glasser}
 {}_kF_{k-1}\!\left(
 \begin{matrix}\frac1{k+1},\frac2{k+1},\ldots,\frac{k}{k+1}\\
 \frac2k,\frac3k,\ldots,\frac{k-1}k,\frac{k+1}k
 \end{matrix};
 \left(\frac{x(1-x^k)}{f_k}\right)^k\right)
 =\frac1{1-x^k}.
\end{equation}
For \(0\le x\le(k+1)^{-1/k}\), both sides have their ordinary positive-real
meaning and the convergence is effective.
\end{theorem}

\begin{proof}
Put \(u=x^k\),
\[
 t=\left(\frac{x(1-x^k)}{f_k}\right)^k,
 \qquad y=\frac1{1-u}.
\]
Since \(f_k^k=k^k/(k+1)^{k+1}\), direct algebra gives
\begin{equation}\label{eq:glasser-algebraic}
 y=1+f_k^k t\,y^{k+1}.
\end{equation}
The coefficient of the unique formal solution is the Fuss--Catalan number
\begin{equation}\label{eq:fuss}
 [t^n]y=\frac1{kn+1}\binom{(k+1)n}{n}f_k^{kn}.
\end{equation}
A fully finitary proof of \eqref{eq:fuss} is obtained by induction from
\eqref{eq:glasser-algebraic}: the coefficient of order \(n\) is a bounded
\((k+1)\)-fold convolution.  Alternatively, the standard Lagrange coefficient
formula is proved once by the finite binomial theorem and coefficient
extraction.

Apply the factorial multiplication identity
\[
 \prod_{j=1}^{k}\poch{j/(k+1)}n
 =\frac{((k+1)n)!}{(k+1)^{(k+1)n}n!}.
\]
After separating multiples of \(k+1\), and the analogous denominator
identity for the fractions with denominator \(k\).  Cross-multiplication
reduces the hypergeometric coefficient exactly to \eqref{eq:fuss}.  Hence the
left side is the unique formal solution of \eqref{eq:glasser-algebraic}, and
substitution gives \(y=(1-x^k)^{-1}\).  On the stated interval, the argument
lies in the convergence range and every root is the positive isolated root.
\end{proof}

\begin{corollary}[The \(k=2\) trigonometric-looking case]
For \(0\le x\le1/\sqrt3\),
\[
 \F{1/3}{2/3}{3/2}{\frac{27}{4}x^2(1-x^2)^2}
 =\frac1{1-x^2}.
\]
Combining this with \cref{thm:arcsin-special} or the cubic equation yields the
intermediate sine expression displayed by MathWorld; the algebraic equality
is the foundationally simpler statement.
\end{corollary}

\subsection{Jacobi, Legendre, and Chebyshev polynomials}

\begin{theorem}[Jacobi polynomial representation]\label{thm:jacobi}
For \(n\in\N\),
\begin{equation}\label{eq:jacobi}
 \F{-n}{n+\alpha+\beta+1}{\alpha+1}{x}
 =\frac{n!}{\poch{\alpha+1}n}P_n^{(\alpha,\beta)}(1-2x),
\end{equation}
where the right side may be characterized as the polynomial solution of the
Jacobi differential equation normalized by
\(P_n^{(\alpha,\beta)}(1)=\poch{\alpha+1}n/n!\).
\end{theorem}

\begin{proof}
The left side terminates at degree \(n\).  Substitute \(z=(1-t)/2\) into the
hypergeometric ODE.  Clearing denominators yields
\[
 (1-t^2)y''+[\beta-\alpha-(\alpha+\beta+2)t]y'
 +n(n+\alpha+\beta+1)y=0,
\]
the Jacobi equation.  The value at \(t=1\) is \(1\), so multiplication by
\(\poch{\alpha+1}n/n!\) gives the standard normalization.  Among polynomials
of degree at most \(n\), the coefficient recurrence determines the solution.
Everything is finite and remains valid over localized rational parameter
rings.
\end{proof}

\begin{corollary}[Legendre polynomials]
\[
 P_n(t)=\F{-n}{n+1}{1}{\frac{1-t}{2}}.
\]
\end{corollary}

\begin{theorem}[Generalized Chebyshev/cosine identity]\label{thm:chebyshev}
Formally in \(s\),
\[
 \F a{-a}{1/2}{\frac{1-s}{2}}=T_a(s).
\]
On a named real interval \(-\pi<x<\pi\), with \(s=\cos x\),
\[
 \F a{-a}{1/2}{\sin^2(x/2)}=\cos(ax).
\]
For integral \(a=n\), this is the ordinary Chebyshev polynomial identity and
is entirely finite.
\end{theorem}

\begin{proof}
Let \(z=\sin^2(x/2)\).  The chain rule transforms the hypergeometric equation
with parameters \((a,-a,1/2)\) into
\(Y''+a^2Y=0\).  The normalized solution has
\(Y(0)=1\), \(Y'(0)=0\), which characterizes \(\cos(ax)\) by its primitive
recursive Taylor coefficients and tail bound.  For integral \(n\), termination
makes \(T_n\) a polynomial and no trigonometric analysis is required.
\end{proof}

\begin{remark}[Legendre functions of nonintegral degree]
The source page also gives a formula involving an associated Legendre
function and gamma factors.  We treat it as a \StatusC{} identity: after a
specific Legendre branch is defined by its differential equation and local
normalization, the equality follows from the same change of variable and
formal uniqueness.  A global statement across its cuts additionally needs
branch-continuation data.
\end{remark}

\section{Transformation formulas}\label{sec:transformations}

\subsection{A finite pullback--gauge certificate}

\begin{lemma}[Pullback certificate]\label{lem:pullback-certificate}
Let \(v(w)\) satisfy \(\mathcal L_{A,B,C}v=0\).  Let
\(w=\phi(z)\) and \(g(z)=p(z)v(\phi(z))\), where \(p,\phi\) are rational or
algebraic named functions.  After differentiating twice and replacing
\(v''\) with the expression supplied by its differential equation, one gets
\[
 \mathcal L_{a,b,c}g=R_0(z)v(\phi(z))+R_1(z)v'(\phi(z)).
\]
If exact algebraic certificates prove \(R_0=R_1=0\), and if
\(g(0)=1\), then
\[
 g(z)=\F a b c z
\]
as a formal series.  On any branch domain with effective convergence, the
identity holds between named analytic functions.
\end{lemma}

\begin{proof}
The chain rule is a finite symbolic calculation.  Clearing all denominators
turns \(R_0=R_1=0\) into polynomial identities in the parameter and pullback
codes.  The resulting \(g\) satisfies the target ODE.  Its constant term fixes
the normalized regular solution by \cref{cor:ode-unique}.  If a resonant
parameter makes the regular recurrence singular, the same certificate is
read after localization or at each finite truncation where the required
denominators are guarded.
\end{proof}

\begin{remark}
This lemma is a practical proof checker for transformation formulas.  It does
not rely on recognizing two analytic functions after the fact.  The
certificate consists only of two cleared rational-function numerators and a
normalization value.
\end{remark}

\subsection{Euler and Pfaff transformations}

\begin{theorem}[Euler transformation]\label{thm:euler-transform}
On the normalized branch at \(z=0\),
\begin{equation}\label{eq:euler-transform}
 \F a b c z=(1-z)^{c-a-b}\F{c-a}{c-b}{c}{z}.
\end{equation}
\end{theorem}

\begin{proof}
Apply \cref{lem:pullback-certificate} with \(\phi(z)=z\),
\(p(z)=(1-z)^{c-a-b}\), and source parameters
\((c-a,c-b,c)\).  Using
\(p'/p=-(c-a-b)/(1-z)\), direct collection makes both residual
coefficients zero.  The prefactor and inner series have constant term \(1\).
The calculation is polynomial after multiplication by \((1-z)^2\).
\end{proof}

\begin{theorem}[Pfaff transformations]\label{thm:pfaff}
On the normalized branch at zero,
\begin{align}
 \F a b c z
 &=(1-z)^{-a}
   \F a{c-b}c{\frac{z}{z-1}},
   \label{eq:pfaff-a}\\
 \F a b c z
 &=(1-z)^{-b}
   \F{c-a}b c{\frac{z}{z-1}}.
   \label{eq:pfaff-b}
\end{align}
\end{theorem}

\begin{proof}
For \eqref{eq:pfaff-a}, use
\(\phi=z/(z-1)\), \(p=(1-z)^{-a}\), and parameters
\((a,c-b,c)\) in \cref{lem:pullback-certificate}.  The derivative data
\[
 \phi'=-\frac1{(z-1)^2},\qquad
 \phi''=\frac2{(z-1)^3},\qquad
 \frac{p'}p=\frac a{1-z}
\]
make the residual numerators vanish.  The second formula follows by symmetry
in \(a,b\).  All denominators are powers of \(1-z\), so the formal identity
is checked in a localized polynomial ring and the analytic identity needs a
branch disk avoiding \(1\).
\end{proof}

\begin{corollary}[Half-argument conversion]\label{cor:half-minus-one}
\[
 \F a b c{1/2}=2^a\F a{c-b}c{-1}
\]
whenever both named values and the chosen power of two are defined.
\end{corollary}

\begin{proof}
Put \(z=1/2\) in \eqref{eq:pfaff-a}.
\end{proof}

\subsection{Quadratic transformations}

\begin{theorem}[A standard quadratic transformation]\label{thm:quadratic}
On the branch normalized at zero,
\begin{equation}\label{eq:quadratic}
 \F a b{2b}{z}
 =(1-z)^{-a/2}
 \F{a/2}{b-a/2}{b+1/2}{\frac{z^2}{4z-4}}.
\end{equation}
\end{theorem}

\begin{proof}
Apply \cref{lem:pullback-certificate} with
\[
 \phi(z)=\frac{z^2}{4z-4},
 \qquad p(z)=(1-z)^{-a/2}.
\]
Use a source equation with parameters
\((a/2,b-a/2,b+1/2)\), and target parameters \((a,b,2b)\).
Multiplication by \(16(1-z)^4\) clears every derivative denominator.  The
coefficient of \(v'\) and then that of \(v\) expand to zero.  At zero,
\(p(0)=1\), \(\phi(0)=0\).  Moreover, the right derivative is
\(a/2=ab/(2b)\), so the normalization is compatible even in a formulation
using the first two coefficients.  This provides a finite checker rather
than an appeal to a table of transformations.
\end{proof}

\begin{theorem}[Kummer quadratic pair]\label{thm:kummer-quadratic}
For the normalized branches,
\begin{align}
 \F a b{1+a-b}{z}
 &=(1-z)^{-a}
 \F{a/2}{(a+1)/2-b}{1+a-b}
   {-\frac{4z}{(1-z)^2}},
 \label{eq:kummer-q1}\\
 \F a b{1+a-b}{z}
 &=(1+z)^{-a}
 \F{a/2}{(a+1)/2}{1+a-b}
   {\frac{4z}{(1+z)^2}}.
 \label{eq:kummer-q2}
\end{align}
\end{theorem}

\begin{proof}
Both are applications of \cref{lem:pullback-certificate}.  For the first use
\(\phi=-4z/(1-z)^2\), \(p=(1-z)^{-a}\); for the second use
\(\phi=4z/(1+z)^2\), \(p=(1+z)^{-a}\).  Clearing \((1\mp z)^4\) reduces the
claims to polynomial identities.  The constant term is \(1\).  The source
parameter relation \(c=1+a-b\) is exactly what cancels the residual
coefficient.
\end{proof}

\begin{remark}[Constructive domain statement]
The formal transformations do not need inequalities.  Their analytic reading
requires that the chosen disk avoid the poles of the pullback and that the
prefactor branch be named there.  The equality may then be continued along a
coded path by \cref{thm:path-continuation}.  This is more precise than a
single informal phrase such as ``by analytic continuation.''
\end{remark}

\subsection{The six local expressions and Kummer's 24 forms}

Let the six base expressions be
\[
 u_{0,1},u_{0,2},u_{1,1},u_{1,2},u_{\infty,a},u_{\infty,b}
\]
from \cref{prop:zero-solutions,prop:one-solutions,prop:infty-solutions}.
For any hypergeometric factor in one of these expressions, apply one of four
operations:
\begin{enumerate}[label=(\arabic*)]
\item leave it unchanged;
\item apply Euler's transformation;
\item apply the first Pfaff transformation;
\item apply the second Pfaff transformation.
\end{enumerate}
After simplifying the outer monomial, this produces the 24 forms tabulated by
Kummer and displayed by MathWorld.

\begin{theorem}[Kummer-form generator]\label{thm:kummer-24}
For generic nonresonant parameters, every expression generated by the
preceding \(6\times4\) construction solves the same hypergeometric differential
equation.  Within the local branch domain of its parent base expression, it
represents that same local solution.  The construction and its verification
are primitive recursive symbolic procedures.
\end{theorem}

\begin{proof}
Each base expression solves the target equation by the operator calculations
of \cref{sec:ode}.  Euler and Pfaff transformations replace only the inner
hypergeometric factor by a formally equal expression, so the parent function
is unchanged in its normalized local branch.  Syntactic simplification of
powers adds exponents and composes rational pullbacks; this is finite rational
arithmetic.  Alternatively, each of the 24 outputs can be checked directly
by \cref{lem:pullback-certificate}.  Resonant coincidences reduce the number
of distinct forms but do not invalidate the equation-solving claim.
\end{proof}

\begin{audit}[What ``all 24'' means]
The 24 entries are not 24 globally single-valued functions.  They are local
expressions containing powers such as \((-z)^{-a}\) and
\((1-z)^{c-a-b}\).  Equality of two entries away from the normalized local
disk includes a connection constant determined by continuation.  The formal
certificate proves the differential equation and local normalization; the
analytic branch statement carries a path and logarithm certificate.
\end{audit}

\subsection{Higher rational pullbacks}

\begin{proposition}[Finite checkability of algebraic transformations]
Suppose a proposed quadratic, cubic, or higher transformation has the form
\[
 \F a b c z=p(z)\F A B C{\phi(z)},
\]
where \(p\) and \(\phi\) are algebraic functions given by isolated polynomial
relations and local initial values.  If implicit differentiation supplies
named \(p',p'',\phi',\phi''\), then the identity has a finite
pullback--gauge certificate as in \cref{lem:pullback-certificate}.
\end{proposition}

\begin{proof}
Use the polynomial relations to eliminate all algebraic powers in the two
residual coefficients.  A Gröbner-style reduction is not foundationally
needed: a certificate explicitly expresses each cleared residual numerator as
a polynomial combination of the defining relations.  Checking that
combination is primitive recursive polynomial arithmetic.  The initial value
selects the intended local branch.
\end{proof}

This proposition covers the specific transformations stated on the source
pages and gives a direct formalization route for the many additional Goursat
transformations that those pages mention but do not print.

\section{Euler integrals and summation theorems}\label{sec:integrals-sums}

\subsection{The beta integral as effective data}

For named complex parameters with positive real-part margins, define
\[
 B(x,y)=\int_0^1 t^{x-1}(1-t)^{y-1}\dd t.
\]
On the positive real axis, the powers use the real logarithm, so
\[
 \abs{t^{x-1}(1-t)^{y-1}}
 =t^{\RePart x-1}(1-t)^{\RePart y-1}.
\]
If \(\RePart x,\RePart y\ge\delta>0\), endpoint tails are bounded by
\(\eta^\delta/\delta\), up to an explicit branch-independent constant.
Thus the cutoff \(\eta\), central quadrature mesh, and evaluation precision
are primitive recursive functions of the requested accuracy.

\begin{lemma}[Beta shift identities]\label{lem:beta-shift}
Under positive real-part margins,
\begin{align}
 B(x+1,y)&=\frac{x}{x+y}B(x,y),\label{eq:beta-shift-x}\\
 B(x,y+1)&=\frac{y}{x+y}B(x,y),\label{eq:beta-shift-y}\\
 B(x+n,y)&=B(x,y)\frac{\poch x n}{\poch{x+y}n}.
 \label{eq:beta-shift-n}
\end{align}
These are effective equalities of names.
\end{lemma}

\begin{proof}
The identity
\[
 1=t+(1-t)
\]
gives \(B(x,y)=B(x+1,y)+B(x,y+1)\).  Integration by parts, performed first
on \([\eta,1-\eta]\), gives
\(xB(x,y+1)=yB(x+1,y)\).  The endpoint terms tend effectively to zero by the
positive margins.  Solving the two linear equations yields
\eqref{eq:beta-shift-x} and \eqref{eq:beta-shift-y}; induction gives
\eqref{eq:beta-shift-n}.  The finite-quadrature proof of integration by parts
uses telescoping sums and a supplied derivative modulus, as in
\cref{lem:finite-substitution}.
\end{proof}

The gamma package in the repository proves an effective version of
\[
 B(x,y)=\frac{\Gamma(x)\Gamma(y)}{\Gamma(x+y)}
\]
in the positive chamber and extends it by the gamma recurrence.  The present
arguments may take either the beta integral or that gamma quotient as the
primary name.

\subsection{Euler's hypergeometric integral}

\begin{theorem}[Euler integral]\label{thm:euler-integral}
Assume rational margins
\[
 \RePart b\ge\delta,
 \qquad
 \RePart(c-b)\ge\delta,
 \qquad \delta>0,
\]
and suppose \(1-zt\) stays apart from zero and from the chosen logarithm cut
for \(0\le t\le1\).  Then
\begin{equation}\label{eq:euler-integral}
 B(b,c-b)\F a b c z
 =\int_0^1t^{b-1}(1-t)^{c-b-1}(1-zt)^{-a}\dd t.
\end{equation}
For \(\norm z\le r<1\), the proof supplies a primitive recursive error
schedule uniform in \(z\).
\end{theorem}

\begin{proof}
On \(\norm z\le r<1\), use the effective binomial expansion
\[
 (1-zt)^{-a}=\sum_{n<N}\frac{\poch a n}{n!}(zt)^n+R_N(t,z),
\]
with a uniform geometric remainder being available from
\cref{thm:disk-convergence}.  Integrating the finite head gives
\[
 \sum_{n<N}\frac{\poch a n}{n!}z^nB(b+n,c-b).
\]
By \eqref{eq:beta-shift-n}, this is
\[
 B(b,c-b)\sum_{n<N}
 \frac{\poch a n\poch b n}{\poch c n n!}z^n.
\]
The integral of the remainder is bounded by the uniform remainder times an
effective bound for the beta weight.  Endpoint cutoffs and the central mesh
are chosen first; then \(N\) is chosen so that the integrated remainder is at
most the remaining error budget.  This proves equality of names.

For a more general \(z\) satisfying the stated branch and apartness condition,
one may either estimate the same integral directly with a local polynomial
approximation to \((1-zt)^{-a}\), or continue the already proved identity
along a certified path.  The latter is \cref{thm:path-continuation}.
\end{proof}

\begin{corollary}[The \([0,\infty)\) Euler form]\label{cor:euler-infty}
Under the corresponding margins and branch data,
\[
 \F a b c z
 =\frac1{B(b,c-b)}
 \int_0^\infty
 u^{b-1}(1+u)^{a-c}(1+(1-z)u)^{-a}\dd u.
\]
\end{corollary}

\begin{proof}
In \eqref{eq:euler-integral}, put \(t=u/(1+u)\).  The finite substitution
principle applies after cutting the improper integral at a primitive
recursive bound and estimating the algebraic tail.
\end{proof}

\subsection{Incomplete beta and elliptic integrals}

\begin{corollary}[Incomplete beta]\label{cor:incomplete-beta}
For positive real-part and branch margins,
\begin{equation}\label{eq:incomplete-beta}
 B_x(p,q)=\int_0^x t^{p-1}(1-t)^{q-1}\dd t
 =\frac{x^p}{p}\F p{1-q}{p+1}{x}.
\end{equation}
\end{corollary}

\begin{proof}
Substitute \(t=xu\).  The remaining integral is Euler's integral with
\(a=1-q\), \(b=p\), \(c=p+1\), and
\(B(p,1)=1/p\).  The branch of \(x^p\) is part of the input data.
\end{proof}

\begin{corollary}[Complete elliptic integrals]\label{cor:elliptic}
For \(|k|<1\),
\begin{align}
 K(k)&=\frac\pi2\F{1/2}{1/2}{1}{k^2},\label{eq:elliptic-K}\\
 E(k)&=\frac\pi2\F{-1/2}{1/2}{1}{k^2}.
 \label{eq:elliptic-E}
\end{align}
\end{corollary}

\begin{proof}
Euler's integral with \(b=1/2\), \(c=1\) has beta factor
\(B(1/2,1/2)=\pi\).  The substitution \(t=\sin^2\theta\) gives
\[
 \frac12t^{-1/2}(1-t)^{-1/2}\dd t=\dd\theta.
\]
For \(a=1/2\) and \(a=-1/2\), the remaining factors are respectively
\((1-k^2\sin^2\theta)^{-1/2}\) and its positive square root.  In iPRA one can
instead take the \(t\)-integrals as the definitions and verify the
trigonometric substitution by finite quadrature; no existence theorem for an
arbitrary antiderivative is needed.
\end{proof}

\subsection{Finite Chu--Vandermonde}

\begin{theorem}[Chu--Vandermonde]\label{thm:vandermonde}
For \(m\in\N\), under the finite denominator guards,
\begin{equation}\label{eq:vandermonde}
 \F{-m}{b}{c}{1}=\frac{\poch{c-b}m}{\poch c m}.
\end{equation}
This is a \StatusA{} theorem: no convergence or gamma function is involved.
\end{theorem}

\begin{proof}
Let
\[
 S_m(b,c)=\sum_{n=0}^m(-1)^n\binom mn\frac{\poch b n}{\poch c n}.
\]
We prove the formula by induction on \(m\), uniformly in \(b,c\).  Pascal's
identity gives
\begin{align*}
 S_{m+1}(b,c)
 &=S_m(b,c)-\frac bc S_m(b+1,c+1).
\end{align*}
By the induction hypothesis,
\[
 S_m(b,c)=\frac{\poch{c-b}m}{\poch c m},
 \qquad
 S_m(b+1,c+1)=\frac{\poch{c-b}m}{\poch{c+1}m}.
\]
Since
\(\poch{c+1}m=\poch c m(c+m)/c\),
\begin{align*}
 S_{m+1}(b,c)
 &=\frac{\poch{c-b}m}{\poch c m}
   \left(1-\frac b{c+m}\right)\\
 &=\frac{\poch{c-b}{m+1}}{\poch c{m+1}}.
\end{align*}
After cross-multiplication this is a quantifier-free identity between finite
products and sums.  The base \(m=0\) is immediate.
\end{proof}

\begin{corollary}[MathWorld's Kummer first formula]
For \(n\in\N\),
\[
 \F{1/2+m-k}{-n}{2m+1}{1}
 =\frac{\Gamma(2m+1)\Gamma(m+1/2+k+n)}
        {\Gamma(m+1/2+k)\Gamma(2m+1+n)},
\]
whenever the finite denominators are guarded.  The gamma notation is merely a
rewriting of a quotient of two rising factorials.
\end{corollary}

\subsection{Gauss's summation theorem}

\begin{theorem}[Effective Gauss summation in the Euler chamber]\label{thm:gauss}
Assume positive rational margins for
\(\RePart b\), \(\RePart(c-b)\), and
\(\RePart(c-a-b)\).  Then
\begin{equation}\label{eq:gauss}
 \F a b c1
 =\frac{\Gamma(c)\Gamma(c-a-b)}
        {\Gamma(c-a)\Gamma(c-b)}.
\end{equation}
The equality has a primitive recursive modulus.
\end{theorem}

\begin{proof}
Set \(z=1\) in Euler's integral.  The factor
\((1-t)^{-a}\) combines with the beta weight, giving
\[
 B(b,c-b)\F a b c1=B(b,c-a-b).
\]
Use the beta--gamma identity twice and cancel the apart gamma factors.  The
endpoint margins provide the integration modulus.  Thus the proof is not an
appeal to a limiting value at the boundary; it directly evaluates an
improper integral with an explicit tail estimate.
\end{proof}

\begin{remark}[Beyond the Euler chamber]
The classical formula only requires \(\RePart(c-a-b)>0\), apart from poles.
Starting from a parameter triple in the Euler chamber, contiguous relations
shift \(a,b,c\) by integers, and symmetry interchanges \(a,b\).  A finite
chain, together with nonzero coefficient certificates, proves many additional
instances.  Full generic extension is a \StatusC{} continuation theorem in
parameter space: the path must avoid the gamma poles and the singular
hyperplanes of the recurrence.
\end{remark}

\subsection{Kummer's theorem at \(z=-1\)}

\begin{theorem}[Kummer summation]\label{thm:kummer-sum}
Under the margins needed by the transformed Gauss sum,
\begin{equation}\label{eq:kummer-sum}
 \F a b{1+a-b}{-1}
 =\frac{\Gamma(1+a-b)\Gamma(1+a/2)}
        {\Gamma(1+a)\Gamma(1+a/2-b)}.
\end{equation}
Equivalently,
\[
 \F a b{1+a-b}{-1}
 =2^{-a}\sqrt\pi\,
  \frac{\Gamma(1+a-b)}
       {\Gamma((1+a)/2)\Gamma(1+a/2-b)}.
\]
\end{theorem}

\begin{proof}
Put \(z=-1\) in \eqref{eq:kummer-q1}.  The transformed argument is \(1\),
so Gauss's theorem gives
\[
 2^{-a}\frac{\Gamma(1+a-b)\Gamma(1/2)}
 {\Gamma((1+a)/2)\Gamma(1+a/2-b)}.
\]
The gamma duplication formula
\[
 \Gamma((1+a)/2)\Gamma(1+a/2)
 =2^{-a}\sqrt\pi\,\Gamma(1+a)
\]
turns this into \eqref{eq:kummer-sum}.  The quadratic transformation is a
finite ODE certificate; the remaining analytic content is the effective
Gauss sum and the already developed gamma identity.
\end{proof}

\subsection{Theorems at \(z=1/2\)}

\begin{theorem}[Gauss's second summation theorem]\label{thm:gauss-second}
Under the corresponding margins,
\begin{equation}\label{eq:gauss-second}
 \F a b{(1+a+b)/2}{1/2}
 =\frac{\Gamma(1/2)\Gamma((1+a+b)/2)}
        {\Gamma((1+a)/2)\Gamma((1+b)/2)}.
\end{equation}
\end{theorem}

\begin{proof}
The pullback certificate of \cref{lem:pullback-certificate} verifies
\begin{equation}\label{eq:gauss-second-quad}
 \F a b{(1+a+b)/2}{z}
 =\F{a/2}{b/2}{(1+a+b)/2}{4z(1-z)}
\end{equation}
near zero.  At \(z=1/2\), the right argument is \(1\).  Gauss's theorem
applies because
\[
 \frac{1+a+b}{2}-\frac a2-\frac b2=\frac12,
\]
and gives exactly \eqref{eq:gauss-second}.  The path from zero to one half
stays in the ordinary convergence disk, so no extra branch choice is needed.
\end{proof}

\begin{theorem}[Bailey's theorem]\label{thm:bailey}
In an initial Euler chamber, and elsewhere by certified parameter
continuation,
\begin{equation}\label{eq:bailey}
 \F a{1-a}c{1/2}
 =\frac{\Gamma(c/2)\Gamma((c+1)/2)}
        {\Gamma((c+a)/2)\Gamma((c+1-a)/2)}.
\end{equation}
\end{theorem}

\begin{proof}
Use Euler's integral with the second numerator parameter \(b=1-a\):
\[
 B(1-a,c+a-1)F
 =\int_0^1t^{-a}(1-t)^{c+a-2}(1-t/2)^{-a}\dd t.
\]
Put \(u=1-t\).  Since
\(t(1-t/2)=(1-u^2)/2\), the integral becomes
\[
 2^a\int_0^1u^{c+a-2}(1-u^2)^{-a}\dd u.
\]
Then put \(v=u^2\), obtaining
\[
 2^{a-1}B((c+a-1)/2,1-a).
\]
Divide by the original beta factor and rewrite with gamma functions:
\[
 F=2^{a-1}
 \frac{\Gamma(c)\Gamma((c+a-1)/2)}
      {\Gamma(c+a-1)\Gamma((c+1-a)/2)}.
\]
Apply the duplication formula once to \(\Gamma(c+a-1)\) and once to
\(\Gamma(c)\).  The powers of two and \(\sqrt\pi\) cancel, leaving
\eqref{eq:bailey}.  Every substitution is covered by
\cref{lem:finite-substitution}; all endpoint exponents have supplied positive
margins.
\end{proof}


\subsection{A further MathWorld value at \(z=-1\)}

\begin{theorem}[The \(\F1{-a}a{-1}\) evaluation]\label{thm:minus-one-extra}
In the initial chamber \(\RePart a>1\), and elsewhere by certified parameter
continuation,
\begin{equation}\label{eq:minus-one-extra}
 \F1{-a}a{-1}
 =1+\frac{\sqrt\pi}{2}\frac{\Gamma(a+1)}{\Gamma(a+1/2)}.
\end{equation}
\end{theorem}

\begin{proof}
Use symmetry and Euler's integral with numerator parameters \((-a,1)\):
\[
 F=(a-1)\int_0^1(1-t)^{a-2}(1+t)^a\dd t.
\]
Put \(t=2u-1\).  Then
\[
 F=(a-1)2^{2a-1}\int_{1/2}^1u^a(1-u)^{a-2}\dd u.
\]
Let \(J\) denote this last integral and let \(K\) be the corresponding
integral from \(0\) to \(1/2\).  Clearly
\(J+K=B(a+1,a-1)\).  Substituting \(u\mapsto1-u\) in \(J\) gives
\[
 J-K=\int_0^{1/2}[u(1-u)]^{a-2}(1-2u)\dd u
 =\frac{4^{1-a}}{a-1},
\]
because the integrand is
\((a-1)^{-1}\frac{d}{du}[u(1-u)]^{a-1}\).  Hence
\[
 J=\frac12\left(B(a+1,a-1)+\frac{4^{1-a}}{a-1}\right).
\]
After multiplication by \((a-1)2^{2a-1}\), the second term is \(1\), while
the first is
\[
 2^{2a-2}\frac{\Gamma(a+1)\Gamma(a)}{\Gamma(2a)}
 =\frac{\sqrt\pi}{2}\frac{\Gamma(a+1)}{\Gamma(a+1/2)}
\]
by duplication.  Every integral substitution has an effective endpoint
margin in the stated chamber.
\end{proof}

\subsection{The isolated MathWorld values}

MathWorld lists the following evaluations:
\begin{align*}
 \F{1/3}{2/3}{5/6}{27/32}&=\frac85,\\
 \F{1/4}{1/2}{3/4}{80/81}&=\frac95,\\
 \F{1/8}{3/8}{1/2}{2400/2401}&=\frac{2\sqrt7}{3},\\
 \F{1/6}{1/3}{1/2}{25/27}&=\frac{3\sqrt3}{4},\\
 \F{1/6}{1/2}{2/3}{125/128}&=\frac43\,2^{1/6},\\
 \F{1/12}{5/12}{1/2}{1323/1331}&=\frac34\,11^{1/4},
\end{align*}
and a further value at \(121/125\) involving
\(\Gamma(1/3)^3/\Gamma(1/4)^2\).

\begin{audit}[Status of the isolated values]\label{audit:isolated-values}
The archive numerically validates the displayed values to high precision and
represents every radical target by an isolated algebraic name.  A completed
iPRA proof of each algebraic value would consist of either:
\begin{enumerate}
\item a finite algebraic pullback certificate reducing the value to a
normalized point already covered above; or
\item a modular-equation certificate, an AGM/modular transformation chain,
and explicit convergence bounds.
\end{enumerate}
The source pages attribute these values to work on special hypergeometric
values, but do not print the transformation certificates.  This package does
not reconstruct all of those modular proofs, so the six displayed isolated
values and the gamma quotient are marked \StatusD{}, not silently counted as
proved.  This is the largest conspicuous omission from the finite list of
formulas on the two pages.
\end{audit}

\section{Contiguous relations and Gauss's continued fraction}\label{sec:contiguous}

Write \(F=\F a b c z\) and \(\theta=zD\).  A symbol such as \(F(a+)\)
means that the indicated parameter is increased by one.

\subsection{The basic contiguous identities}

\begin{theorem}[Seven expressions for \(\theta F\)]\label{thm:contiguous-master}
Under the relevant finite denominator guards,
\begin{align}
 \theta F
 &=z\frac{ab}{c}F(a+,b+,c+),\label{eq:contig-der}\\
 &=a(F(a+)-F),\label{eq:contig-a+}\\
 &=b(F(b+)-F),\label{eq:contig-b+}\\
 &=(c-1)(F(c-)-F),\label{eq:contig-c-}\\
 &=\frac{(c-a)F(a-)+(a-c+bz)F}{1-z},\label{eq:contig-a-}\\
 &=\frac{(c-b)F(b-)+(b-c+az)F}{1-z},\label{eq:contig-b-}\\
 &=z\frac{(c-a)(c-b)F(c+)+c(a+b-c)F}{c(1-z)}.
 \label{eq:contig-c+}
\end{align}
Every equality is a formal power-series identity and hence a \StatusA{}
iPRA theorem schema.
\end{theorem}

\begin{proof}
Equation \eqref{eq:contig-der} is the first derivative formula multiplied by
\(z\).  For \eqref{eq:contig-a+}, compare coefficients:
\[
 \frac{\poch{a+1}n}{\poch a n}=\frac{a+n}{a},
 \qquad
 a\left(\frac{a+n}{a}-1\right)A_n=nA_n.
\]
The cross-multiplied calculation remains valid at \(a=0\).  The \(b\) formula
is symmetric.  Similarly,
\[
 \frac{\poch c n}{\poch{c-1}n}=\frac{c+n-1}{c-1}
\]
gives \eqref{eq:contig-c-}.

For \eqref{eq:contig-a-}, multiply by \(1-z\) and compare the coefficient of
\(z^n\).  The coefficient from \(zF\) is the preceding coefficient
\(A_{n-1}\), which is converted to \(A_n\) using
\eqref{eq:ratio}.  Clearing
\(n(n+c-1)\) reduces the assertion to a polynomial identity in
\(n,a,b,c\).  The \(b\) identity follows by symmetry.  The last identity is
proved in the same way using
\(\poch{c+1}n/\poch c n=(c+n)/c\).  Thus all seven lines have direct bounded
coefficient checkers.
\end{proof}

\begin{corollary}[The fifteen Gauss relations]
Equating any two of the seven right-hand expressions in
\cref{thm:contiguous-master} yields the usual relations expressing \(F\) as a
linear combination of any two among its six contiguous functions.  After
clearing the displayed denominators, every relation is a polynomial identity
of formal series.
\end{corollary}

\subsection{Arbitrary integer shifts}

\begin{theorem}[Constructive contiguous reduction]\label{thm:shift-reduction}
For any integer triples \(r_i=(m_i,n_i,\ell_i)\), \(i=1,2,3\), there are
rational functions \(R_1,R_2,R_3\), not all zero, such that
\[
 \sum_{i=1}^3R_i(a,b,c,z)
 \F{a+m_i}{b+n_i}{c+\ell_i}{z}=0,
\]
relative to explicit guards for the finitely many denominators introduced by
the reduction.  The coefficient lists of \(R_i\) are primitive recursive in
the shift triples.
\end{theorem}

\begin{proof}
Use \cref{thm:contiguous-master} to move one parameter at a time toward a
chosen pair of basis functions, for example \(F\) and \(F(a+)\).  A measure
\(|m|+|n|+|\ell|\) decreases at every step.  Each step is a \(2\times2\)
rational matrix; multiply the finite matrices and clear their denominators.
Three vectors in the resulting two-dimensional module are linearly dependent,
and a cross-product gives \(R_1,R_2,R_3\).  The algorithm is a bounded loop
whose length is controlled by the shift measure.  No appeal to the dimension
of an abstract analytic solution space is needed, although the result mirrors
that classical explanation.
\end{proof}

\subsection{Exact finite continued fractions}

Define
\begin{align*}
 f_{2k}(z)&=\F{a+k}{b+k}{c+2k}{z},\\
 f_{2k+1}(z)&=\F{a+k+1}{b+k}{c+2k+1}{z},
\end{align*}
and
\begin{align*}
 \alpha_{2k+1}
 &=\frac{(a-c-k)(b+k)}{(c+2k)(c+2k+1)},\\
 \alpha_{2k+2}
 &=\frac{(b-c-k-1)(a+k+1)}{(c+2k+1)(c+2k+2)}.
\end{align*}

\begin{lemma}[Gauss difference chain]\label{lem:gauss-chain}
For every \(j\ge1\),
\begin{equation}\label{eq:gauss-chain}
 f_{j-1}-f_j=\alpha_j z f_{j+1}.
\end{equation}
\end{lemma}

\begin{proof}
For odd and even \(j\), substitute the corresponding shifted parameters into
a basic contiguous relation.  Coefficientwise, both cases reduce to the
Pochhammer successor identity.  The checker needs only the parity function,
the half function, and rational arithmetic, all available by primitive
recursion.
\end{proof}

Let \(r_j=f_j/f_{j-1}\), relative to apartness certificates for the
denominators.  Dividing \eqref{eq:gauss-chain} by \(f_{j-1}\) gives
\begin{equation}\label{eq:ratio-tail}
 r_j=\frac1{1+\alpha_jzr_{j+1}}.
\end{equation}

\begin{theorem}[Finite Gauss fraction with an exact tail]\label{thm:finite-cf}
For every \(N\ge1\),
\begin{equation}\label{eq:finite-cf}
 \frac{\F{a+1}{b}{c+1}{z}}{\F a b c z}
 =\cfrac1{1+\cfrac{\alpha_1z}{1+\cfrac{\alpha_2z}{1+
 \cfrac{\ddots}{1+\alpha_Nzr_{N+1}}}}}.
\end{equation}
This is an exact finite rational identity of named functions.
\end{theorem}

\begin{proof}
Iterate \eqref{eq:ratio-tail} exactly \(N\) times.  Formally one may avoid all
divisions by using continuants.  Define numerator and denominator polynomials
of the finite fraction by the two-term recurrence; induction shows that their
cross-product difference with \(f_1/f_0\) is zero after the terminal tail
\(f_{N+1}/f_N\) is inserted.
\end{proof}

\subsection{The infinite continued fraction}

Replacing the exact terminal tail \(r_{N+1}\) by \(1\) gives the usual
\(N\)-th convergent printed on the source page.

\begin{contract}[Continued-fraction convergence certificate]\label{con:cf}
A convergence certificate on a rational domain \(D\) consists of:
\begin{enumerate}
\item apartness margins for every \(f_j\) used up to the requested depth;
\item a primitive recursive index \(J\) and rational \(\rho<1\) such that the
Möbius maps \(T_j(w)=(1+\alpha_jzw)^{-1}\) map a fixed rational disk \(\Delta\)
into itself and are \(\rho\)-Lipschitz there for all \(j\ge J\), uniformly in
\(z\in D\);
\item a proof that both the exact tail \(r_j\) and the substituted tail \(1\)
lie in \(\Delta\).
\end{enumerate}
\end{contract}

\begin{theorem}[Certified Gauss continued fraction]\label{thm:gauss-cf}
Under \cref{con:cf}, the convergents obtained from
\eqref{eq:finite-cf} by putting the terminal tail equal to \(1\) converge
uniformly and effectively to
\[
 \frac{\F{a+1}{b}{c+1}{z}}{\F a b c z}.
\]
A valid depth for precision \(2^{-p}\) is
\(J+m\), where bounded search finds \(m\) with
\(L\rho^m\le2^{-p}\) and \(L\) bounds the finite prefix.
\end{theorem}

\begin{proof}
The exact and substituted tails are propagated through the same composition
\(T_1\circ\cdots\circ T_N\).  Beyond \(J\), each additional map contracts
the discrepancy by \(\rho\).  The finite prefix has a computable Lipschitz
constant because its denominators have apartness margins.  Thus the final
error is at most \(L\rho^{N-J}\).  The proof is induction on a finite
composition of rational maps.
\end{proof}

\begin{audit}[Why the infinite fraction is \StatusC]
The finite identity \cref{thm:finite-cf} is unconditional apart from finitely
many guards and is fully algebraic.  Convergence of the infinite fraction is
not obtained merely by printing an ellipsis.  Zeros of denominator functions
and parameter-dependent stability matter.  The certificate
\cref{con:cf} isolates exactly the extra analytic content.  On many compact
subdisks it can be discharged from explicit coefficient bounds, but a single
uniform proof covering every parameter regime displayed classically is not
claimed here.
\end{audit}

\section{Effective continuation, connection formulas, and global claims}\label{sec:continuation}

\subsection{Taylor stepping along a coded path}

A qualitative statement that a function ``can be analytically continued along
any path avoiding the singularities'' must be replaced by a finite path and
quantitative separation data.

\begin{definition}[Continuation certificate]\label{def:continuation-cert}
A continuation certificate from \(z_0\) to \(z_m\) consists of rational or
complex-rational centers
\(z_0,z_1,\ldots,z_m\), positive rational radii \(R_j\), and margins
\(d_j>0\) such that
\begin{align*}
 \norm{z_{j+1}-z_j}&\le R_j/2,\\
 \operatorname{dist}(\overline{D(z_j,R_j)},\{0,1\})&\ge d_j,
\end{align*}
plus initial names for \(y(z_0),y'(z_0)\).  If powers occur, the certificate
also contains compatible logarithm branches on the disks and overlap
constants.
\end{definition}

On a disk avoiding \(0,1\), rewrite the equation as
\[
 y''=P(z)y'+Q(z)y,
 \quad
 P(z)=-\frac{c-(a+b+1)z}{z(1-z)},
 \quad
 Q(z)=\frac{ab}{z(1-z)}.
\]
Expand \(P,Q\) at the center.  Their coefficients have geometric bounds from
the separation margin.  The Taylor coefficients of \(y\) then satisfy a
finite convolution recurrence determined by \(y(z_j),y'(z_j)\).

\begin{theorem}[Primitive recursive path continuation]\label{thm:path-continuation}
Given \cref{def:continuation-cert}, a primitive recursive algorithm computes
names for \(y(z_j),y'(z_j)\) at every center and hence a name at the endpoint.
The result is independent, up to effective equality, of refinements of the
same certified disk chain.  Applied to the regular series at zero after an
initial center inside \(|z|<1\), it gives an effective branch of
\(\F a b c z\) along the path.
\end{theorem}

\begin{proof}
For one step, majorize the coefficient recurrence by a scalar series
\(M/(1-w/R)\) whose rational constants are computed from
\(d_j,R_j,\norm a,\norm b,\norm c\) and the current value bounds.  Induction
shows that the Taylor coefficient of order \(n\) is at most \(MR^{-n}\).
At a point with displacement at most \(R/2\), the tail after \(N\) terms is
at most \(2M2^{-N}\), and the differentiated tail has a similarly explicit
bound.  Choose \(N\) by bounded search, evaluate the finite Taylor polynomial,
and propagate the resulting intervals to the next center.  Repeat for the
finite list of disks.

If a disk chain is refined, both constructions solve the same local initial
value problem on each overlap.  The coefficient recurrence proves uniqueness,
and the two tail moduli transfer the coefficient equality to equality of
names.  This is an induction over the finite refinement, not a compactness
argument over all paths.
\end{proof}

\begin{corollary}[Branch cut convention]
A conventional principal branch on
\(\C\setminus[1,\infty)\) is represented by a family of rational polygonal
paths from a base point, each carrying a positive distance from the cut and
from \(0,1\).  Equality of two evaluations is proved when their paths are
connected by an explicitly coded homotopy grid whose cells stay a positive
distance from the singularities.
\end{corollary}

\subsection{The connection formula at one}

For nonresonant parameters, define
\begin{align*}
 A&=\frac{\Gamma(c)\Gamma(c-a-b)}
         {\Gamma(c-a)\Gamma(c-b)},\\
 B&=\frac{\Gamma(c)\Gamma(a+b-c)}
         {\Gamma(a)\Gamma(b)}.
\end{align*}
The classical connection formula is
\begin{equation}\label{eq:connection-one}
\begin{split}
 \F a b c z={}&A\,
 \F a b{a+b+1-c}{1-z}\\
 &+B(1-z)^{c-a-b}
 \F{c-a}{c-b}{1+c-a-b}{1-z}.
\end{split}
\end{equation}

\begin{theorem}[Certified connection formula]\label{thm:connection}
Equation \eqref{eq:connection-one} is a \StatusC{} iPRA theorem schema under:
\begin{enumerate}
\item apartness from the gamma poles and from the resonant hyperplanes
\(c-a-b\in\Z\);
\item a branch of \((1-z)^{c-a-b}\);
\item a finite parameter-continuation chain from a chamber in which the
connection coefficients are determined by effective Gauss evaluations;
\item a path-continuation certificate in the \(z\)-plane.
\end{enumerate}
\end{theorem}

\begin{proof}
Both terms on the right solve the original differential equation by
\cref{prop:one-solutions}.  Thus it suffices to determine the two constants in
one nonresonant parameter chamber.  One coefficient is the regular limit at
\(z=1\) and equals \(A\) by \cref{thm:gauss}.  Applying Euler's
transformation exchanges the two exponents at one; in a chamber where the
transformed regular limit exists, Gauss's theorem gives \(B\).  To place both
calculations in one formal proof, choose a finite path in parameter space that
avoids poles and resonance, and use contiguous relations whenever an integer
shift is more convenient.  Gamma recurrence and the ODE uniqueness theorem
show that the constants continue along the same chain.  Finally continue in
\(z\) with \cref{thm:path-continuation}.

Every step has finite data.  What is conditional is the supplied path and its
margins; no unrestricted theorem that the complement of all singular
hyperplanes is path connected is used inside iPRA.
\end{proof}

\begin{remark}[Resonant connection formulas]
When \(c-a-b\in\Z\), gamma coefficients in
\eqref{eq:connection-one} have poles and the two Frobenius solutions coalesce.
Taking an effective parameter limit produces a logarithmic formula involving
finite harmonic sums and gamma derivatives.  This can be formalized for any
fixed resonance and truncation order, but a uniform library of all resonant
cases is not included.
\end{remark}

\subsection{Barnes's contour integral}

The Barnes integrand is
\[
 G(s)=
 \frac{\Gamma(a+s)\Gamma(b+s)\Gamma(-s)}{\Gamma(c+s)}
 (-z)^s.
\]
A vertical contour must separate the poles
\(s=0,1,2,\ldots\) from
\(s=-a-n,-b-n\).  This separation is not expressible by a picture alone.

\begin{contract}[Barnes contour data]\label{con:barnes}
Supply:
\begin{enumerate}
\item a rational vertical line \(\RePart s=\sigma\), or a finite polygonal
deformation, with positive distances from every relevant pole family;
\item a branch of \(\log(-z)\) with
\(|\ImPart\log(-z)|\le\pi-\eta\) for rational \(\eta>0\);
\item explicit vertical-strip gamma bounds of the form
\[
 \norm{G(\sigma+it)}\le C(1+|t|)^M e^{-\eta|t|};
\]
\item for a closing rectangle, an explicit bound on the far vertical side,
typically containing \(\norm z^N\) when \(\norm z<1\).
\end{enumerate}
The gamma/Stirling and Cauchy-residue packets in the repository provide the
intended lower-level ingredients.
\end{contract}

\begin{theorem}[Certified Barnes integral]\label{thm:barnes}
Under \cref{con:barnes},
\begin{equation}\label{eq:barnes}
 \frac1{2\pi i}\int_{-i\infty}^{i\infty}G(s)\dd s
 =\frac{\Gamma(a)\Gamma(b)}{\Gamma(c)}\F a b c z
\end{equation}
with the contour interpreted as the supplied separating path.  The equality
has a primitive recursive precision schedule.
\end{theorem}

\begin{proof}
First truncate the contour at height \(T\) and close it after the first
\(N\) poles of \(\Gamma(-s)\).  The finite residue theorem gives a finite sum
of residues.  At \(s=n\), the residue, including the orientation sign, is
\[
 \frac{\Gamma(a+n)\Gamma(b+n)}{\Gamma(c+n)n!}z^n
 =\frac{\Gamma(a)\Gamma(b)}{\Gamma(c)}
   \frac{\poch a n\poch b n}{\poch c n n!}z^n.
\]
This is an exact gamma-recurrence calculation.  The horizontal edges and the
infinite vertical tails are bounded by
\(C(1+T)^Me^{-\eta T}\); bounded search chooses \(T\).  The far vertical edge
and the uncollected series tail are bounded by the supplied \(N\)-estimate;
choose \(N\) next.  The total error is allocated among the finitely many
edges and the series tail.  Thus the ordinary contour proof becomes a finite
residue identity plus primitive recursive inequalities.
\end{proof}

\begin{audit}
The report specifies the complete proof architecture but does not contain a
kernel-checked implementation of all complex gamma strip bounds and contour
integration.  Accordingly Barnes's formula is counted as \StatusC{}, not as
an unconditional finished analytic theorem.
\end{audit}

\subsection{Local monodromy and Schwarz ratios}

For a nonresonant Frobenius basis at a singular point with exponents
\(\rho_1,\rho_2\), continuation once around a small certified polygonal loop
multiplies the basis by
\[
 \begin{pmatrix}e^{2\pi i\rho_1}&0\\0&e^{2\pi i\rho_2}\end{pmatrix}.
\]
The statement is effective for rational or algebraically named exponents,
because the exponential has a Taylor modulus and the loop is finite.

\begin{proposition}[Möbius action on a solution ratio]\label{prop:schwarz-local}
Let \(s=\varphi_1/\varphi_0\) on a disk where \(\varphi_0\) is apart from
zero.  If continuation changes the basis by
\(\binom{\varphi_0}{\varphi_1}\mapsto
 M\binom{\varphi_0}{\varphi_1}\), then
\[
 s\longmapsto\frac{M_{10}+M_{11}s}{M_{00}+M_{01}s}.
\]
Moreover,
\[
 s'=\frac{W(\varphi_0,\varphi_1)}{\varphi_0^2},
\]
so local conformality is certified by apartness of the Wronskian and
\(\varphi_0\).
\end{proposition}

\begin{proof}
Both identities are one-line cross-multiplications.  The Wronskian has the
explicit form of \cref{lem:wronskian}, so its nonvanishing on a branch disk is
quantitative once the leading constant is apart from zero.
\end{proof}

\begin{audit}[Global Schwarz and monodromy claims]
The source page describes triangle maps, explicit monodromy matrices, and the
classification saying that certain rational exponent differences give finite
monodromy exactly when \(1/k+1/l+1/m>1\).  The local matrices and their finite
products are arithmetic objects and can be checked once connection matrices
are certified.  The global conformal mapping theorem and the complete finite
group classification are not reconstructed here.  They remain \StatusD{}.
\end{audit}

\subsection{The modular-lambda statement}

From \cref{cor:elliptic},
\[
 \F{1/2}{1/2}1z=\frac2\pi K(\sqrt z).
\]
Hence the ratio printed on the source page is
\[
 \tau=i\frac{K(\sqrt{1-z})}{K(\sqrt z)}
\]
for compatible square-root branches.  The assertion
\(\lambda(\tau)=z\), where lambda is a quotient of theta series, is the
classical elliptic-modular inversion theorem.

\begin{audit}[Modular inversion]
Theta series at an explicitly named \(\tau\) have elementary Gaussian tail
moduli, so their values are representable in iPRA.  What is missing is a
complete constructive proof that the theta quotient inverts the elliptic
integral ratio for every certified domain point, including the modular
transformation and global branch bookkeeping.  The package therefore proves
the reduction to the elliptic-modular theorem but marks the inversion itself
\StatusD{}.  The same missing modular infrastructure explains the status of
several isolated MathWorld values in \cref{audit:isolated-values}.
\end{audit}

\subsection{The John-transform sentence}

The Wikipedia page states that the Gauss function can be written as a John
transform but does not display a formula or normalization.  A mathematical
statement cannot be formalized from a bare reference to a transform with
multiple conventions.  Once an explicit kernel, domain, and normalization
are supplied, the intended proof is expected to be a finite change of
variables reducing the transform to Euler's integral.  No such theorem is
counted in the coverage total.

\section{Lean formalization plan}\label{sec:formalization}

\subsection{Two distinct Lean layers}

A faithful formalization should not confuse a theorem in Lean's ambient
classical mathematics with a derivation in the object theory \(\ipra\).
We therefore propose two linked libraries.

\begin{enumerate}
\item \textbf{Mathlib/model layer.}  Define ordinary finite coefficients and
analytic sums over \(\Q,\R,\C\); prove the algebraic identities, convergence
bounds, and correctness of certificate generators.  This layer is useful for
rapid development and for checking formulas.
\item \textbf{Object-theory layer.}  Represent every primitive recursive
coefficient, partial sum, cutoff, and error test by a symbol/term in
\code{Prf PRAAx}; derive the first-order sentences using the checked iPRA core.
Only the second layer establishes the foundational claim of this report.
\end{enumerate}

The repository already contains a Mathlib-side definition
\code{Fabius.classicalHypergeometricTerm} for \({}_{s+1}F_s\), using
Mathlib's rising-factorial polynomial, and a noncomputable infinite sum
\code{classicalHypergeometric}.  It also proves nonvanishing and termination
lemmas in the context of a classical \(q\to1\) limit.  These definitions can
be reused for the model layer, but the noncomputable \(\sum'\) is not itself
an iPRA construction.

The checked \code{IPRADerivations} library supplies the object-theory
infrastructure needed here: integer differences, a ring tactic, definitions
by primitive recursion, finite sums and products, powers, factorials, order,
bounded search, pairing, and multi-sums.  This sharply reduces the amount of
new foundational code.

\subsection{Proposed module graph}

\begin{center}
\small
\begin{tabularx}{\textwidth}{@{}p{.28\textwidth}X@{}}
\toprule
Module & Main declarations \\
\midrule
\path{Hypergeometric/Codes} & signed rationals, complex-rational pairs,
normalized or cross-multiplied equality, one-norm, dyadic error predicate \\
\path{Hypergeometric/Pochhammer} & rising factorial numerator/denominator,
split and shift laws, negative-integer termination \\
\path{Hypergeometric/Coefficient} & \(A_n\), guard predicate, ratio
recurrence, generalized \({}_pF_q\) term \\
\path{Hypergeometric/FiniteSeries} & partial sums, Cauchy products,
derivatives, substitution in \(R/(z^N)\) \\
\path{Hypergeometric/ODE} & coefficient form of \(\mathcal L\), normalized
uniqueness, Frobenius gauge identities, Q-form \\
\path{Hypergeometric/Transform} & Euler, two Pfaff, quadratic/Kummer
certificates, six/24-form generator \\
\path{Hypergeometric/Special} & logarithm, arcsine, cubic equation, Glasser
family, Jacobi/Legendre/Chebyshev \\
\path{Hypergeometric/Bounds} & explicit ratio cutoff, geometric tail,
head--tail combinators, derivative moduli \\
\path{Hypergeometric/Integral} & endpoint cutoffs, Riemann sums, beta shifts,
Euler integral, substitutions \\
\path{Hypergeometric/Summation} & Vandermonde, Gauss chamber, Kummer,
Gauss second, Bailey \\
\path{Hypergeometric/Contiguous} & seven master identities, shift reducer,
finite Gauss fraction \\
\path{Hypergeometric/Continuation} & local Taylor recurrence, disk-step
majorant, finite path continuation \\
\path{Hypergeometric/Checks/Axioms} & \code{\#print axioms} for every exported
object-theory theorem; no \code{sorry} \\
\bottomrule
\end{tabularx}
\end{center}

The exact finite core should be implemented first.  It is independent of the
analytic-name infrastructure and will already formalize a substantial fraction
of both reference pages.

\subsection{Executable model definitions}

The companion file \code{lean/HypergeometricFinite.lean} uses a direct
rational definition suitable for experimentation:
\begin{lstlisting}[language={}]
def pochhammer (a : Q) : Nat -> Q
| 0     => 1
| n + 1 => pochhammer a n * (a + n)

def hgCoeff (a b c : Q) (n : Nat) : Q :=
  pochhammer a n * pochhammer b n /
    (pochhammer c n * n.factorial)

def hgPartial (a b c z : Q) (N : Nat) : Q :=
  sumRange N (fun n => hgCoeff a b c n * z^n)
\end{lstlisting}
The production version should use a guarded rational structure or state
nonzero hypotheses, rather than relying on Lean's total field division at a
zero denominator.

At the Mathlib level, specialization of the existing generic definition is:
\begin{lstlisting}[language={}]
def gaussTerm (a b c z : K) (n : Nat) : K :=
  Fabius.classicalHypergeometricTerm
    (s := 1) (fun | 0 => a | 1 => b) (fun _ => c) z n
\end{lstlisting}
One then proves agreement with \code{hgCoeff} over \(\Q\) and uses the
existing lemmas about \code{ascPochhammer}.

\subsection{The object-theory representation}

Working directly with rational quotients makes cancellation side conditions
awkward.  A robust object-theory implementation stores a rational expression
as a pair of signed-integer terms \((N,D)\), with \(D>0\).  Equality is
\[
  N_1D_2=N_2D_1
\]
in the ring of integer differences.  Addition and multiplication are
polynomial; inversion swaps numerator and denominator under a nonzero guard.
For most hypergeometric identities it is even cheaper to keep the coefficient
as the unreduced pair
\[
 \operatorname{num}(A_n)=\poch a n\poch b n,
 \qquad
 \operatorname{den}(A_n)=\poch c n\,n!,
\]
where each rational parameter has first been expanded into a common integer
denominator.  The core recurrence is then one integer polynomial equality.

A typical target theorem has the shape
\begin{lstlisting}[language={}]
theorem coeff_recurrence :
  ([] : List (Fm PTm 0)) |- [PRAAx]
    allN 4 (guardCoeff ==>.
      ratEq
        ((n+1) * (n+c) * coeff (n+1))
        ((n+a) * (n+b) * coeff n)) := by
  -- unfold primitive recursions; use product equations and ipra_ring
\end{lstlisting}
This is schematic notation; the archive's Lean outline records the intended
binder order and imports, not a claim that these exact identifiers already
exist.

\subsection{Induction invariants for the finite core}

The crucial proofs have unusually small invariants.

\paragraph{Pochhammer split.}
For fixed \(x,m\), induct on \(n\) with the literal equation
\[
 \poch x{m+n}=\poch x m\poch{x+m}n.
\]

\paragraph{Termination.}
Induct on \(n\le m\) for the factorial formula, then use the zero factor at
\(m+1\) and the successor equation for later indices.

\paragraph{Derivative order.}
Induct on \(r\) with a universally quantified observed coefficient \(n\), or
encode the bounded conjunction of all coefficients below \(N\).  The step is
\(\poch a{r+1}=\poch a r(a+r)\).

\paragraph{ODE.}
No induction is needed beyond the coefficient definitions: one
\code{ipra\_ring} invocation proves the cross-multiplied recurrence at an
arbitrary \(n\).

\paragraph{Vandermonde.}
Use the induction in \cref{thm:vandermonde}.  The induction hypothesis is
universal in the rational parameter codes, so it may be instantiated at
\((b+1,c+1)\) in the successor step.

\paragraph{Transformations.}
Implement a small symbolic differential-expression datatype.  A certificate
contains two rational-polynomial residuals and identities expressing them as
zero.  Reflection into the object theory turns each residual check into a
ring theorem.  Formal uniqueness then transports the result to coefficients.

\subsection{Formalizing the analytic contracts}

A dyadic complex-rational approximation can be stored as four natural
numbers: signed numerators for real and imaginary parts and a common power-of
two denominator.  Define
\[
 \operatorname{errLE}(X,Y,p)
\]
to mean, after scaling to a common denominator,
\(\norm{X-Y}\le2^{-p}\).  This is a quantifier-free formula over natural
arithmetic.

For a fixed evaluator \(E(p)\), the Cauchy contract is the sentence
\[
 \forall p\forall q\;
 (q\ge p+s\to\operatorname{errLE}(E(q),E(p+s),p)).
\]
The analytic equality of fixed evaluators \(E_1,E_2\) is
\[
 \forall p\;
 \operatorname{errLE}(E_1(p+s),E_2(p+s),p).
\]
The quantifiers are ordinary first-order quantifiers; only the evaluator
symbols are fixed by the theorem schema.

The tail theorem can be compiled without a general theory of real limits.
Define the primitive recursive cutoff \(N(a,b,c,r,p)\) by the explicit formula
and bounded search of \cref{lem:ratio-cutoff,thm:disk-convergence}.  Prove
inside iPRA that every finite tail segment between \(N\) and \(M\) is bounded
by the geometric expression.  The Cauchy theorem follows directly.  This
approach is shorter than first building a full abstract constructive analysis
library.

\subsection{Formalizing Euler's integral}

A minimal development needs only the following reusable lemmas:
\begin{enumerate}
\item a bounded-sum estimate for rational complex numbers;
\item a modulus-to-Riemann-sum Cauchy lemma on a closed rational interval;
\item endpoint power-tail bounds with a positive rational exponent margin;
\item summation/integration interchange for a finite head plus a uniform
remainder;
\item finite substitution for the specific rational maps used in this report.
\end{enumerate}
Each theorem is a first-order statement about explicit sums and inequalities.
The integral is not introduced as a higher-order object; every named integrand
gets its own primitive recursive Riemann-sum symbol.

\subsection{Suggested implementation milestones}

\begin{enumerate}[label=\textbf{M\arabic*.}]
\item Compile Pochhammer, coefficients, termination, derivatives, ODE,
formal uniqueness, and Vandermonde.  This should contain no analytic names.
\item Add a finite truncated-series evaluator and reflection checker.  Prove
Euler, Pfaff, the printed quadratic transformations, the cubic equation, and
the Glasser family.
\item Add rational/complex-rational bounds and prove the strict-disk tail
modulus and termwise differentiation.
\item Add beta Riemann sums and prove the Euler-chamber integral and Gauss
summation; reuse the gamma packet for the quotient form.
\item Add Kummer, Gauss second, Bailey, incomplete beta, and elliptic
specializations.
\item Add path stepping and conditional certificates for the connection
formula, continued fraction, and Barnes contour.
\end{enumerate}

\subsection{Expected theorem declarations}

The final object library should export, at minimum, declarations corresponding
to the following mathematical statements:
\begin{lstlisting}[language={}]
poch_split, poch_shift, poch_neg, poch_zero_after
hg_coeff_ratio, hg_terminates, hg_deriv_coeff
hg_ode_coeff, hg_regular_unique
hg_euler_formal, hg_pfaff1_formal, hg_pfaff2_formal
hg_quadratic_formal, hg_kummer_quadratic_formal
hg_contiguous_master, hg_shift_reduction
chu_vandermonde
hg_cubic_algebraic, glasser_family
hg_ratio_cutoff, hg_tail_cauchy, hg_termwise_derivative
euler_integral_effective, gauss_sum_effective
kummer_sum_effective, gauss_second_effective, bailey_effective
hg_path_continue, gauss_fraction_finite
\end{lstlisting}
The last two analytic groups should state all margins and moduli in their
types; hidden typeclass assumptions such as abstract completeness should not
be used to discharge them.

\subsection{Audit status of the supplied Lean files}

The archive includes:
\begin{itemize}
\item \code{HypergeometricFinite.lean}, a detailed Mathlib-style finite-core
outline with definitions and many proof scripts or proof sketches;
\item \code{HypergeometricIPRA.lean}, the object-theory declaration plan and
formula encodings;
\item \code{FORMALIZATION\_PLAN.md}, a file-by-file dependency and proof
checklist.
\end{itemize}
The environment used to prepare this package did not contain a Lean executable,
so these files were not compiled.  They are explicitly labeled as an outline;
no claim of kernel verification is made.  The existing repository files cited
above are separate and already provide checked infrastructure.

\section{Coverage ledger and limitations}\label{sec:inventory}

\subsection{How the ledger is counted}

The companion file \code{theorem\_inventory.csv} has one row for each
mathematical assertion or displayed formula family identified on the two
reference pages.  Closely related assertions appearing on both pages are kept
as separate rows, so the ledger measures source coverage rather than the
number of mathematically independent theorems.  There are 104 rows:
\begin{center}
\begin{tabular}{@{}lrr@{}}
\toprule
Status & Rows & Interpretation \\
\midrule
\StatusA & 16 & complete finite/formal iPRA proof \\
\textsf{A/B} & 32 & complete finite proof and effective analytic reading \\
\textsf{A/C} & 11 & complete local/formal proof; branch/global part certified \\
\StatusB & 20 & effective analytic proof with explicit moduli \\
\StatusC & 14 & conditional proof with a displayed finite certificate contract \\
\StatusD & 11 & not proved in this package \\
\midrule
Total & 104 & \\
\bottomrule
\end{tabular}
\end{center}
Thus 48 rows have a complete exact finite core (statuses \StatusA{} and
\textsf{A/B}), another 20 are effective analytic theorems, and 25 have either
a complete local theorem or a precise finite certificate reduction.  Eleven
rows remain genuinely outside the present reduction.

\subsection{What is strongest in the package}

The broadest exact result is not any one special value but the reusable
pullback--gauge checker of \cref{lem:pullback-certificate}.  It reduces a
large class of transformations to two polynomial residuals.  In the same
spirit, \cref{thm:contiguous-master,thm:shift-reduction} turn all integer
parameter shifts into a bounded rational matrix computation, and
\cref{thm:finite-cf} proves the finite continued fraction with its exact
hypergeometric tail rather than treating the ellipsis as self-justifying.

The main analytic contribution is the separation of convergence into an
explicit rational cutoff and a geometric remainder
(\cref{lem:ratio-cutoff,thm:disk-convergence}).  It supplies a direct route to
object-theory Cauchy sentences and avoids importing a classical ratio-test
library.  Euler's integral and the major summation theorems then use finite
head--tail and endpoint-cutoff arguments, which are naturally compatible with
iPRA.

Two source formulas benefit from particularly economical constructive proofs:
\begin{itemize}
\item the cubic radical is the isolated branch of
\(Y+x^2Y^3=1\), identified with the hypergeometric series by ODE uniqueness;
\item the entire Glasser family is the Fuss--Catalan solution of
\(y=1+f_k^kty^{k+1}\), with the hypergeometric coefficient obtained by finite
factorial multiplication identities.
\end{itemize}

\subsection{The eleven \StatusD{} rows}

The omissions fall into four groups.

\paragraph{Seven isolated modular values.}
MathWorld's equations (18)--(24) are verified numerically in the archive, but
their published modular or algebraic transformation proofs are not
reconstructed.  Numerical agreement is not counted as proof.

\paragraph{Global modular inversion.}
The lambda/theta inversion statement is reduced to the classical inversion of
the elliptic integral ratio, but a complete effective modular-function
library is absent.

\paragraph{Finite-monodromy classification.}
Local monodromy matrices and finite products are covered conditionally; the
global classification by spherical triangle groups is not.

\paragraph{Under-specified mentions.}
The source page only mentions several other orthogonal polynomial families and
a John-transform representation without printing formulas or normalization.
No theorem is invented to fill that gap.

\subsection{Full source-row table}

The compact table below gives the row identifier, source claim, status, and
report route.  The CSV adds source locators, formalization methods, and notes.

\begingroup\scriptsize
\setlength{\tabcolsep}{3pt}
\begin{longtable}{@{}p{.065\textwidth}p{.53\textwidth}p{.08\textwidth}p{.245\textwidth}@{}}
\toprule ID & Statement & Status & Report route \\
\midrule\endfirsthead
\toprule ID & Statement & Status & Report route \\
\midrule\endhead
\midrule\multicolumn{4}{r}{\emph{continued}}\\\endfoot
\bottomrule\endlastfoot
W01 & Gauss 2F1 power-series definition & \textsf{A/B} & Def. 3.3; Thm. 4.2 \\
W02 & Rising factorial product and recurrences & \textsf{A} & Def. 3.1; Lem. 3.2 \\
W03 & Ordinary coefficients fail when c is a nonpositive integer & \textsf{A} & Def. 3.3 \\
W04 & Termination and explicit polynomial for a=-m & \textsf{A} & Thm. 3.6 \\
W05 & Continuation along paths avoiding 1 and infinity & \textsf{C} & Thm. 10.2 \\
W06 & Limit of 2F1/Gamma(c) as c approaches -m & \textsf{B} & Thm. 3.8 \\
W07 & 2F1 as the p=2,q=1 generalized series & \textsf{A} & Def. 3.9 \\
W08 & First derivative parameter shift & \textsf{A/B} & Thm. 3.12; Thm. 4.4 \\
W09 & Arbitrary finite derivative order & \textsf{A/B} & Thm. 3.12; Thm. 4.4 \\
W10 & 2F1(1,1;2;-z)=log(1+z)/z & \textsf{A/B} & Thm. 6.3 \\
W11 & 2F1(a,b;b;z)=(1-z)\textasciicircum{}(-a) & \textsf{A/B} & Thm. 6.1 \\
W12 & Arcsine hypergeometric representation & \textsf{A/B} & Thm. 6.4 \\
W13 & Cubic-radical representation & \textsf{A/B} & Thm. 6.5 \\
W14 & 2F1(1,b;b;z)=1/(1-z) & \textsf{A/B} & Cor. 6.2 \\
W15 & 1F1 as b->0 limit of 2F1(a,b\textasciicircum{}-1;c;bz) & \textsf{B} & Thm. 4.8 \\
W16 & Functions reducible through confluent/0F1 limits & \textsf{C} & Sec. 4.5; Sec. 11 \\
W17 & Associated Legendre representation & \textsf{C} & Rem. 6.11 \\
W18 & Jacobi polynomial representation & \textsf{A} & Thm. 6.8 \\
W19 & Krawtchouk, Meixner, Meixner-Pollaczek mentions & \textsf{D} & Sec. 12 \\
W20 & Theta quotient inverts a hypergeometric ratio & \textsf{D} & Sec. 10.7 \\
W21 & Incomplete beta hypergeometric identity & \textsf{B} & Cor. 8.4 \\
W22 & Complete elliptic K representation & \textsf{B} & Cor. 8.5 \\
W23 & Complete elliptic E representation & \textsf{B} & Cor. 8.5 \\
W24 & Gauss hypergeometric differential equation & \textsf{A/B} & Thm. 5.1 \\
W25 & Regular singular points at 0,1,infinity & \textsf{A} & Sec. 5.2 \\
W26 & Any certified three-regular-singularity equation reduces to Gauss & \textsf{C} & Thm. 5.9 \\
W27 & Regular solution at zero & \textsf{A/B} & Prop. 5.3 \\
W28 & Second nonresonant solution z\textasciicircum{}(1-c)F(...)  & \textsf{A/C} & Prop. 5.3 \\
W29 & Logarithmic second solution for integral exponent difference & \textsf{C} & Rem. 5.6 \\
W30 & Two local solutions at one & \textsf{A/C} & Prop. 5.4 \\
W31 & Two local solutions at infinity & \textsf{A/C} & Prop. 5.5 \\
W32 & Six Frobenius expressions & \textsf{A/C} & Sec. 5.2 \\
W33 & Twenty-four Kummer forms & \textsf{A/C} & Thm. 7.8 \\
W34 & Euler/Pfaff equivalences among Kummer forms & \textsf{A/B} & Thms. 7.2-7.3 \\
W35 & Elimination of first derivative and printed Q(z) & \textsf{A} & Prop. 5.8 \\
W36 & Ratios of pairs of solutions & \textsf{C} & Prop. 10.7 \\
W37 & Connection matrices act by Mobius transformations & \textsf{A/C} & Prop. 10.7 \\
W38 & Printed monodromy matrices from exponents and connection data & \textsf{C} & Sec. 10.6 \\
W39 & Finite iff 1/k+1/l+1/m>1 in stated rational case & \textsf{D} & Audit 10.8 \\
W40 & Euler beta integral representation & \textsf{B} & Thm. 8.2 \\
W41 & Barnes contour representation & \textsf{C} & Thm. 10.6 \\
W42 & Existence of a John-transform representation & \textsf{D} & Sec. 10.8 \\
W43 & Six contiguous functions and basic identities & \textsf{A} & Thm. 9.1 \\
W44 & Fifteen pairwise contiguous relations & \textsf{A} & Cor. 9.2 \\
W45 & Linear relation among any three integer-shifted functions & \textsf{A} & Thm. 9.3 \\
W46 & Gauss ratio continued fraction & \textsf{A/C} & Thms. 9.5, 9.7 \\
W47 & Euler transformation & \textsf{A/B} & Thm. 7.2 \\
W48 & First Pfaff transformation & \textsf{A/B} & Thm. 7.3 \\
W49 & Second Pfaff transformation & \textsf{A/B} & Thm. 7.3 \\
W50 & Two-term connection formula & \textsf{C} & Thm. 10.4 \\
W51 & Existence pattern for quadratic transformations & \textsf{C} & Sec. 7.3 \\
W52 & Printed transformation for 2F1(a,b;2b;z) & \textsf{A/B} & Thm. 7.5 \\
W53 & Mention of cubic and higher transformations & \textsf{C} & Sec. 7.6 \\
W54 & Gauss summation theorem & \textsf{B} & Thm. 8.8 \\
W55 & Chu-Vandermonde terminating sum & \textsf{A} & Thm. 8.6 \\
W56 & Kummer summation theorem & \textsf{B} & Thm. 8.10 \\
W57 & Two Kummer quadratic transformations & \textsf{A/B} & Thm. 7.6 \\
W58 & Gauss second summation theorem & \textsf{B} & Thm. 8.11 \\
W59 & Bailey summation theorem & \textsf{B} & Thm. 8.12 \\
W60 & Algebraic Chebyshev form with x\textasciicircum{}2/(4(x-1)) & \textsf{A/C} & Thm. 6.10 \\
W61 & 2F1(a,-a;1/2;sin\textasciicircum{}2(x/2))=cos(ax) & \textsf{A/B} & Thm. 6.10 \\
M01 & Generalized hypergeometric successive-term ratio & \textsf{A} & Prop. 3.10 \\
M02 & Hypergeometric ODE & \textsf{A/B} & Thm. 5.1 \\
M03 & Frobenius coefficient recurrence & \textsf{A} & Lem. 3.4; Thm. 5.1 \\
M04 & Regular series solution & \textsf{A/B} & Defs. 3.3-3.5 \\
M05 & Convergence for |z|<1 & \textsf{B} & Thm. 4.2 \\
M06 & Convergence at |z|=1 under positive excess & \textsf{B} & Thm. 4.6 \\
M07 & Two-term complete local solution at zero & \textsf{A/C} & Prop. 5.3 \\
M08 & First derivative formula & \textsf{A/B} & Thm. 3.12; Thm. 4.4 \\
M09 & Second derivative formula & \textsf{A/B} & Thm. 3.12; Thm. 4.4 \\
M10 & 2F1(1,1;1;z)=1/(1-z) & \textsf{A/B} & Cor. 6.2 \\
M11 & 2F1(1,1;2;z)=-log(1-z)/z & \textsf{A/B} & Thm. 6.3 \\
M12 & 2F1(1,2;1;z)=1/(1-z)\textasciicircum{}2 & \textsf{A/B} & Cor. 6.2 \\
M13 & 2F1(1,2;2;z)=1/(1-z) & \textsf{A/B} & Cor. 6.2 \\
M14 & Euler integral & \textsf{B} & Thm. 8.2 \\
M15 & Barnes contour integral & \textsf{C} & Thm. 10.6 \\
M16 & 2F1(1/3,2/3;5/6;27/32)=8/5 & \textsf{D} & Audit 8.14 \\
M17 & 2F1(1/4,1/2;3/4;80/81)=9/5 & \textsf{D} & Audit 8.14 \\
M18 & Value 2 sqrt(7)/3 at 2400/2401 & \textsf{D} & Audit 8.14 \\
M19 & Value 3 sqrt(3)/4 at 25/27 & \textsf{D} & Audit 8.14 \\
M20 & Value (4/3)2\textasciicircum{}(1/6) at 125/128 & \textsf{D} & Audit 8.14 \\
M21 & Value (3/4)11\textasciicircum{}(1/4) at 1323/1331 & \textsf{D} & Audit 8.14 \\
M22 & Gamma quotient at 121/125 & \textsf{D} & Audit 8.14 \\
M23 & Glasser well-poised algebraic family & \textsf{A/B} & Thm. 6.6 \\
M24 & k=2 Glasser specialization and rational value & \textsf{A/B} & Cor. 6.7 \\
M25 & Euler/Pfaff transformations and substitutions & \textsf{A/B} & Thms. 7.2-7.3 \\
M26 & Connection formula at one & \textsf{C} & Thm. 10.4 \\
M27 & Six Kummer solutions & \textsf{A/C} & Props. 5.3-5.5 \\
M28 & Twenty-four Kummer forms & \textsf{A/C} & Thm. 7.8; App. 13 \\
M29 & Existence of Goursat cubic transformations & \textsf{C} & Sec. 7.6 \\
M30 & Legendre polynomial representation & \textsf{A} & Cor. 6.9 \\
M31 & (1+z)\textasciicircum{}n terminating hypergeometric identity & \textsf{A} & Cor. 3.7 \\
M32 & log(1+z)=z 2F1(1,1;2;-z) & \textsf{A/B} & Thm. 6.3 \\
M33 & Complete elliptic integral representations & \textsf{B} & Cor. 8.5 \\
M34 & Kummer z=-1 formula in duplication form & \textsf{B} & Thm. 8.10 \\
M35 & 2F1(1,-a;a;-1) gamma evaluation & \textsf{B} & Thm. 8.13 \\
M36 & z=1/2 to z=-1 Pfaff conversion & \textsf{A/B} & Cor. 7.4 \\
M37 & Gauss second theorem & \textsf{B} & Thm. 8.11 \\
M38 & Bailey theorem & \textsf{B} & Thm. 8.12 \\
M39 & Gauss theorem & \textsf{B} & Thm. 8.8 \\
M40 & Kummer first formula with -n & \textsf{A} & Cor. 8.7 \\
M41 & Generalized nFm definition & \textsf{A/B} & Def. 3.9; Sec. 4 \\
M42 & 1F1 confluent function as a limiting family & \textsf{B} & Thm. 4.8 \\
M43 & 0F1 coefficient definition and strict-disk convergence & \textsf{A/B} & Def. 3.9; Sec. 4 \\
\end{longtable}

\endgroup

\subsection{Claims deliberately weakened or parameterized}

Several classical statements are true in broader domains than the initial
constructive theorem stated here.  The restriction is always visible:
\begin{enumerate}
\item Gauss summation is proved directly in the Euler chamber; broader
instances use a certified parameter-continuation path.
\item The connection formula excludes resonance and asks for a branch/path
certificate.
\item Boundary convergence asks for a positive rational witness
\(\delta\le\RePart(c-a-b)\).
\item Barnes's integral asks for explicit strip and contour-tail bounds.
\item Infinite continued-fraction convergence asks for denominator apartness
and a contraction domain; the finite identity is unconditional.
\item The generic nonresonant Frobenius basis is not asserted in resonant
integer-difference cases; those need logarithmic recurrences.
\end{enumerate}
These are not cosmetic qualifications.  They are the computational data that
turn a classical existence proof into a theorem of primitive recursive
arithmetic.

\subsection{Trust boundary}

The human proofs in this report are intended to be mathematically complete at
the level indicated by their status.  The finite Python checker provides
independent exact tests of coefficient identities and high-precision tests of
selected analytic values.  It is not a proof assistant.  The Lean files are a
formalization outline and were not compiled in the preparation environment.
Only a future derivation in \code{Prf PRAAx}, audited with
\code{\#print axioms} and no \code{sorry}, should be called a completed formal
iPRA proof.

\appendix
\section{The twenty-four Kummer expressions}\label{app:kummer24}

This appendix records a complete working list of the twenty-four classical
expressions generated from one local hypergeometric solution by the six
fractional-linear arguments
\[
 z,\quad 1-z,\quad z^{-1},\quad (1-z)^{-1},\quad
 \frac{z}{z-1},\quad 1-z^{-1}
\]
and the elementary powers of \(z\), \(1-z\), and \(-z\).  They are grouped
into six quartets.  Within a simply connected rational polygonal domain on
which all required powers have certified logarithm branches, each expression
is a solution of the same pulled-back differential equation.  Equality of two
expressions is therefore reduced, as in \cref{thm:formal-uniqueness}, to equality
of two initial coefficients.  This is a \StatusB{} theorem schema.  Without
branch certificates, the table is only a formal list of local expressions;
no global single-valued equality is intended.

For readability write \(F(a,b;c;w)=\F{a}{b}{c}{w}\).  A convenient set of
representatives is
\begin{align*}
K_{1,1}&=F(a,b;c;z),\\
K_{1,2}&=(1-z)^{c-a-b}F(c-a,c-b;c;z),\\
K_{1,3}&=(1-z)^{-a}F\left(a,c-b;c;\frac{z}{z-1}\right),\\
K_{1,4}&=(1-z)^{-b}F\left(c-a,b;c;\frac{z}{z-1}\right);
\end{align*}
\begin{align*}
K_{2,1}&=F(a,b;a+b+1-c;1-z),\\
K_{2,2}&=z^{1-c}F(a+1-c,b+1-c;a+b+1-c;1-z),\\
K_{2,3}&=z^{-a}F(a,a+1-c;a+b+1-c;1-z^{-1}),\\
K_{2,4}&=z^{-b}F(b+1-c,b;a+b+1-c;1-z^{-1});
\end{align*}
\begin{align*}
K_{3,1}&=(-z)^{-a}F(a,a+1-c;a+1-b;z^{-1}),\\
K_{3,2}&=(-z)^{b-c}(1-z)^{c-a-b}
 F(1-b,c-b;a+1-b;z^{-1}),\\
K_{3,3}&=(1-z)^{-a}F\left(a,c-b;a+1-b;(1-z)^{-1}\right),\\
K_{3,4}&=(-z)^{1-c}(1-z)^{c-a-1}
 F\left(a+1-c,1-b;a+1-b;(1-z)^{-1}\right);
\end{align*}
\begin{align*}
K_{4,1}&=(-z)^{-b}F(b+1-c,b;b+1-a;z^{-1}),\\
K_{4,2}&=(-z)^{a-c}(1-z)^{c-a-b}
 F(1-a,c-a;b+1-a;z^{-1}),\\
K_{4,3}&=(1-z)^{-b}F\left(b,c-a;b+1-a;(1-z)^{-1}\right),\\
K_{4,4}&=(-z)^{1-c}(1-z)^{c-b-1}
 F\left(b+1-c,1-a;b+1-a;(1-z)^{-1}\right);
\end{align*}
\begin{align*}
K_{5,1}&=z^{1-c}F(a+1-c,b+1-c;2-c;z),\\
K_{5,2}&=z^{1-c}(1-z)^{c-a-b}F(1-a,1-b;2-c;z),\\
K_{5,3}&=z^{1-c}(1-z)^{c-a-1}
 F\left(a+1-c,1-b;2-c;\frac{z}{z-1}\right),\\
K_{5,4}&=z^{1-c}(1-z)^{c-b-1}
 F\left(b+1-c,1-a;2-c;\frac{z}{z-1}\right);
\end{align*}
\begin{align*}
K_{6,1}&=(1-z)^{c-a-b}F(c-a,c-b;c+1-a-b;1-z),\\
K_{6,2}&=z^{1-c}(1-z)^{c-a-b}
 F(1-a,1-b;c+1-a-b;1-z),\\
K_{6,3}&=z^{a-c}(1-z)^{c-a-b}
 F(c-a,1-a;c+1-a-b;1-z^{-1}),\\
K_{6,4}&=z^{b-c}(1-z)^{c-a-b}
 F(c-b,1-b;c+1-a-b;1-z^{-1}).
\end{align*}

The first quartet consists of the identity, Euler's transformation, and the
two Pfaff transformations.  The remaining quartets are obtained by changing
the distinguished singular point and choosing either local exponent there.
At resonant parameters some displayed pairs coalesce or their normalizations
become singular.  The correct replacement is then a logarithmic Frobenius
solution, so the generic twenty-four-label statement must not be read as
asserting twenty-four linearly independent functions.

\section{Primitive recursive precision schedules}\label{app:schedules}

The arguments in the main text repeatedly use the same small collection of
precision compilers.  This appendix makes them explicit enough to serve as an
implementation specification.

\subsection{Geometric domination}

Let \(0\le r<1\) be rational and let \(A\ge0\).  Put
\[
 d=\operatorname{den}(1-r),\qquad
 E(q,A,r)=\min\left\{N\le B:
 A r^N\le 2^{-q}(1-r)\right\},
\]
where a primitive recursive search bound is
\[
 B=(q+\sigma(A)+2)d,
 \qquad
 \sigma(A)=\min\{s\le \lceil A\rceil+1:2^s\ge\max(1,A)\}.
\]
The elementary inequality \((1-1/d)^d\le 1/2\), proved after clearing
denominators by induction on \(d\), shows that the bounded search succeeds
after enlarging \(B\) by at most a fixed linear factor.  Consequently
\[
 \sum_{n\ge E(q,A,r)} A r^n\le 2^{-q}.
\]
All occurrences of a phrase such as ``choose \(N\) so that the geometric tail
is below \(2^{-q}\)'' can be replaced by this bounded search.

\subsection{Hypergeometric ratio schedule}

For fixed rational-complex parameters with denominator margins
\(\norm{c+n}\ge\eta(n)>0\), the exact ratio is
\[
 R_n(z)=\frac{(a+n)(b+n)}{(c+n)(n+1)}z.
\]
After clearing the parameter denominators, the inequality
\(\norm{R_n(z)}\le r<1\) is a polynomial inequality in \(n\).  The proof of
\cref{lem:ratio-cutoff} supplies an explicit polynomial bound \(N_0\) beyond
which it holds.  The complete cutoff is
\[
 N(q)=\max\bigl(N_0,E(q,A,r)\bigr),
\]
where \(A\) bounds the normalized \(N_0\)-th term.  This is primitive
recursive in the input codes and requested precision.

\subsection{Endpoint truncation for Euler's integral}

Assume rational witnesses
\[
 0<\delta_b\le\RePart b,
 \qquad
 0<\delta_{c-b}\le\RePart(c-b).
\]
For an integrand bounded by
\(M t^{\delta_b-1}(1-t)^{\delta_{c-b}-1}\), choose dyadic
\(\varepsilon=2^{-m}\) by bounded search so that
\[
 \frac{M}{\delta_b}\varepsilon^{\delta_b}
 +\frac{M}{\delta_{c-b}}\varepsilon^{\delta_{c-b}}
 \le 2^{-q-2}.
\]
Rational powers are evaluated by the certified exponential/logarithm names of
\cref{sec:arithmetic}.  On \([\varepsilon,1-\varepsilon]\), use the supplied
modulus of continuity to choose a mesh whose Riemann error is at most
\(2^{-q-2}\), and evaluate each node to precision \(2^{-q-2}/(M+1)\).  The
three errors add to at most \(2^{-q}\).  Every search is bounded by a function
of the displayed rational data.

\subsection{Path continuation schedule}

For a polygonal path \(\gamma\) avoiding \(\{0,1\}\), a continuation
certificate contains rational margins
\[
 \rho_j\le \min_{z\in S_j}\bigl(\abs z,\abs{1-z}\bigr).
\]
After decomposing the path into segments \(S_j\).  On each segment the ODE is
rewritten as a first-order system with a rational Lipschitz bound \(L_j\).
A truncated Taylor or Picard step of length \(h_j\) is accepted when
\(L_jh_j\le 1/4\).  The order is chosen by the geometric compiler above,
and the local errors are budgeted as
\(2^{-q-j-2}\).  Finite summation of the local budgets proves the final
\(2^{-q}\) error.  This realizes continuation without a compactness or
choice principle.

\section{Object-theory target schemas}\label{app:targets}

A theorem in the present report becomes a closed first-order iPRA sentence
only after the primitive recursive programs naming its parameters and
certificates have been fixed.  The following patterns are the intended
translation targets in ProveIt's \code{Prf PRAAx} calculus.

\subsection{Exact coefficient identity}

If \(A(n)\) and \(B(n)\) are integer-coded rational expressions, equality is
represented after cross multiplication by a natural-number equation
\[
 \operatorname{num}A(n)\operatorname{den}B(n)
 =\operatorname{num}B(n)\operatorname{den}A(n).
\]
The target is the universal closure of this quantifier-free equation.  A
simultaneous recurrence can be encoded by the pairing machinery already
available in the repository.

\subsection{Equality of effective complex numbers}

For primitive recursive complex-rational names \(X(q)\) and \(Y(q)\), the
statement \(X=Y\) is compiled to the schema
\[
 \forall q\;
 2^q\,\norm{X(q+2)-Y(q+2)}\le 1.
\]
The norm is the rational \(\ell^1\)-norm, so the displayed predicate is a
finite integer inequality.  The two lost bits absorb the independent naming
errors.  Uniform identities add the bounded parameter code as another free
number variable.

\subsection{Convergence with a modulus}

For partial sums \(S(N)\) and a primitive recursive modulus \(M(q)\), Cauchy
convergence is the arithmetical formula
\[
 \forall q\,\forall m,n\;
 \bigl(M(q)\le m\wedge M(q)\le n\bigr)
 \to 2^q\norm{S(m)-S(n)}\le1.
\]
When the proof uses a geometric tail, the implication is discharged by the
exact finite geometric identity and the cutoff of \cref{app:schedules}.

\subsection{Apartness and branches}

A guarded division \(x/y\) carries an integer witness \(k\) satisfying
\(2^{-k}\le\abs y\).  A branch certificate for \(z^\alpha\) is a finite list
of rational polygons, overlap points, logarithm names, and estimates asserting
that the image stays a rational positive distance from the selected cut.  All
checks are bounded conjunctions of integer inequalities.  Thus the apparent
higher-type data are replaced by one number coding a finite certificate and a
primitive recursive checker.

\subsection{Suggested closed milestone theorems}

The first object-theory implementation should target, in order:
\begin{enumerate}
\item the Pochhammer recurrence and zero criterion;
\item the coefficient ratio, differentiation identity, and ODE residual;
\item terminating Chu--Vandermonde and Jacobi/Legendre specializations;
\item the three basic contiguous relations and their finite continued-fraction
      eliminant;
\item the cubic radical coefficient recurrence;
\item the generic disk tail theorem for a fixed rational parameter tuple;
\item Euler's integral for one fixed positive rational parameter tuple;
\item Gauss summation in the same tuple, followed by a parameterized schema.
\end{enumerate}
The first five use only finite algebra.  The last three exercise the effective
analysis interfaces without yet requiring branch continuation.

\section{Archive map and reproducibility}\label{app:archive}

\begin{center}
\begingroup\small
\begin{tabularx}{\textwidth}{@{}p{0.35\textwidth}X@{}}
\toprule
Path & Role\\
\midrule
\path{report.tex}, \path{report.pdf} & Complete paper and compiled rendering.\\
\path{sections/} & Source sections, including the generated inventory table.\\
\path{theorem_inventory.csv} & Source-row coverage ledger with status, route, and qualifications.\\
\path{lean/HypergeometricFinite.lean} & Executable-looking finite algebra layer and theorem statements; uncompiled outline.\\
\path{lean/HypergeometricIPRA.lean} & Object-theory translation skeleton aimed at \code{Prf PRAAx}.\\
\path{lean/FORMALIZATION_PLAN.md} & Dependency graph, induction invariants, and completion checklist.\\
\path{code/validate_finite.py} & Exact rational and high-precision numerical regression checks.\\
\path{evidence/validation.json} & Deterministic output of the validation program.\\
\path{SOURCE_MANIFEST.md} & Web and repository sources, scope, and inspected revision.\\
\path{SHA256SUMS.txt} & Checksums for all packaged files except the checksum file itself.\\
\bottomrule
\end{tabularx}
\endgroup
\end{center}

The PDF is rebuilt with \code{make pdf}; the finite evidence with
\code{make validate}; and both plus checksums with \code{make all}.  The
validation program proves no general theorem, but it is useful for detecting
sign, index-shift, and normalization mistakes in the human and Lean plans.


\begin{thebibliography}{99}
\bibitem{wiki}
Wikipedia contributors, ``Hypergeometric function,'' revision consulted
16 November 2026. \url{https://en.wikipedia.org/wiki/Hypergeometric_function}.

\bibitem{mathworld}
E. W. Weisstein, ``Hypergeometric Function,'' MathWorld--A Wolfram Web
Resource, consulted 16 November 2026.
\url{https://mathworld.wolfram.com/HypergeometricFunction.html}.

\bibitem{dlmf}
A. B. Olde Daalhuis, ``Hypergeometric Function,'' in NIST Digital Library of
Mathematical Functions, Chapter 15. \url{https://dlmf.nist.gov/15}.

\bibitem{bailey}
W. N. Bailey, \emph{Generalized Hypergeometric Series}, Cambridge University
Press, 1935.

\bibitem{erdelyi}
A. Erd\'elyi, W. Magnus, F. Oberhettinger, and F. G. Tricomi,
\emph{Higher Transcendental Functions}, Vol. I, McGraw--Hill, 1953.

\bibitem{andrews}
G. E. Andrews, R. Askey, and R. Roy, \emph{Special Functions}, Cambridge
University Press, 1999.

\bibitem{proveit}
V. Reshetnikov, \emph{ProveIt}, repository snapshot of the main branch
consulted 16 November 2026.
\url{https://github.com/VladimirReshetnikov/ProveIt}.

\bibitem{gamma-packet}
\emph{Gamma in intuitionistic primitive recursive arithmetic}, in
\cite{proveit}, directory
\path{Logic/PrimitiveRecursiveArithmetic/docs/proofs/analysis/gamma_ipra/},
15 November 2026.

\bibitem{taylor-cauchy}
\emph{Taylor's theorem and Cauchy's integral formula in iPRA}, in
\cite{proveit}, directory
\path{Logic/PrimitiveRecursiveArithmetic/docs/proofs/analysis/taylor_cauchy_ipra/},
15 November 2026.

\bibitem{zucker}
I. J. Zucker and G. S. Joyce, papers on special values of hypergeometric
series cited by MathWorld; the present package records their displayed
values but does not reconstruct the modular-function proofs of all of them.
\end{thebibliography}

\end{document}
